\PassOptionsToPackage{main=italian}{babel}
\documentclass[11pt]{scrbook}
\usepackage{evan}
\usepackage[utf8]{inputenc}
\usepackage[italian]{babel}
\usepackage{algorithm2e}
\usepackage{amsfonts}
\usepackage{amsthm}
\usepackage{amssymb}
\usepackage{amsopn}
\usepackage{biblatex}
\usepackage{bm}
\usepackage{cancel}
\usepackage{csquotes}
\usepackage{mathtools}
\usepackage{marvosym}
\usepackage{tikz}
\usepackage{wrapfig}
\usepackage{xpatch}

\addbibresource{bibliography.bib}

\begin{document}

\title{L'Algebrario}
\subtitle{dispense del corso di Aritmetica}
\author{Gabriel Antonio Videtta}
\date{A.A. 2022/2023 \\ \vskip 1in \includegraphics[scale=0.3]{logo.png}}
\maketitle

\newpage
\thispagestyle{empty}
~\newpage

\section*{Premessa}

TODO

\newpage
\thispagestyle{empty}
~\newpage

\tableofcontents

\newpage
\thispagestyle{empty}
~\newpage

\PassOptionsToPackage{main=italian}{babel}
\documentclass[11pt]{scrbook}
\usepackage{evan_notes}
\usepackage[utf8]{inputenc}
\usepackage[italian]{babel}
\usepackage{amsfonts}
\usepackage{amsthm}
\usepackage{amssymb}
\usepackage{amsopn}
\usepackage{biblatex}
\usepackage{bm}
\usepackage{csquotes}
\usepackage{mathtools}
\usepackage{marvosym}

\begin{document}

\chapter{Introduzione alla teoria degli anelli}

\section{Definizione e prime proprietà}

\begin{definition}
    Si definisce \textbf{anello}\footnote{In realtà, si parla in questo caso di anello \textit{con unità}, in cui vale l'assioma di esistenza di un'identità
        moltiplicativa. In queste dispense si identificherà con "anello" solamente un anello con unità.} una struttura algebrica
    costruita su un insieme $A$ e due operazioni binarie $+$
    e $\cdot$\footnote{D'ora in avanti il punto verrà omesso.} avente le seguenti proprietà:

    \begin{itemize}
        \item $\left(A,\, +\right)$ è un \textit{gruppo abeliano}, alla cui
              identità, detta \textit{identità additiva}, ci si riferisce con il simbolo $0$,
        \item $\forall a, b, c \in A$, $(ab)c = a(bc)$,
        \item $\forall a, b, c \in A$, $(a+b)c=ac+bc$,
        \item $\forall a, b, c \in A$, $a(b+c)=ab+ac$,
        \item $\exists 1 \in A \mid \forall a \in A$, $1a=a=a1$, e tale $1$ viene
              detto \textit{identità moltiplicativa}.
    \end{itemize}
\end{definition}

Come accade per i gruppi, gli anelli soddisfano alcune proprietà algebriche
particolari, tra le quali si citano le più importanti:

\begin{proposition}
    $\forall a \in A$, $0a=0=a0$.
\end{proposition}

\begin{proof}
    $0a=(0+0)a=0a+0a \implies 0a=0$. Analogamente $a0=a(0+0)=a0+a0 \implies a0=0$.
\end{proof}

\begin{proposition}
    $\forall a \in A$, $-(-a)=a$.
\end{proposition}

\begin{proof}
    $-(-a)-a=0 \,\land\, a-a=0 \implies -(-a)=a$, per la proprietà di unicità
    dell'inverso in un gruppo\footnote{In questo caso, il gruppo additivo dell'anello.}.
\end{proof}

\begin{proposition}
    \label{prop:inverso_inverso}
    $a(-b)=(-a)b=-(ab)$.
\end{proposition}

\begin{proof}
    $a(-b)+ab=a(b-b)=a0=0 \implies a(-b)=-(ab)$, per la proprietà di unicità dell'inverso in un gruppo. Analogamente $(-a)b+ab=(a-a)b=0b=0 \implies
        (-a)b=-(ab)$.
\end{proof}

\begin{corollary}
    $(-1)a=a(-1)=-a$.
\end{corollary}

\begin{proposition}
    $(-a)(-b)=ab$.
\end{proposition}

\begin{proof}
    $(-a)(-b)=-(a(-b))=-(-(ab))=ab$, per la \textit{Proposizione \ref{prop:inverso_inverso}}.
\end{proof}

Si enuncia invece adesso la nozione di \textbf{sottoanello}, in tutto e per
tutto analoga a quella di \textit{sottogruppo}.

\begin{definition}
    Si definisce sottoanello rispetto all'anello $A$ un anello $B$ avente le
    seguenti proprietà:

    \begin{itemize}
        \item $B \subseteq A$,
        \item $0, 1 \in B$,
        \item $\forall a, b \in B,$ $a + b \in B \,\land\, ab \in B$.
    \end{itemize}
\end{definition}

\begin{definition}
    Un sottoanello $B$ rispetto ad $A$ si dice \textbf{proprio} se
    $B \neq A$.
\end{definition}

\begin{definition}
    Un anello si dice \textbf{commutativo} se $\forall a$, $b \in A$, $ab=ba$.
\end{definition}

\begin{example}
    Un facile esempio di anello commutativo è $\ZZ/n\ZZ$.
\end{example}

\begin{definition}
    Un elemento $a$ di un anello $A$ si dice \textbf{invertibile} se
    $\exists b \in A \mid ab = ba = 1$.
\end{definition}

\begin{definition}
    Dato un anello $A$, si definisce $A^*$ come l'insieme degli elementi
    invertibili di $A$, che a sua volta forma un \textit{gruppo moltiplicativo}.
\end{definition}

\begin{definition}
    Un anello $A$ si dice \textbf{corpo} se $\forall a \neq 0 \in A$, $\exists b \in A \mid ab=ba=1$,
    ossia se $A \setminus \{0\} = A^*$.
\end{definition}

\begin{example}
    L'esempio più rilevante di corpo è quello dei \textit{quaternioni} $\HH$, definiti
    nel seguente modo:

    \[\HH = \{a+b\ii+c\jj+d\kk \mid a,\, b,\, c,\, d \in \RR\},\]

    dove:

    \[\ii^2 = \jj^2 = \kk^2 = -1, \quad \ii\jj = \kk,\, \jj\kk = \ii,\, \kk\ii = \jj. \]

    Infatti ogni elemento non nullo di $\HH$ possiede un inverso moltiplicativo:

    \[\left(a+b\ii+c\jj+d\kk\right)^{-1} = \frac{a-b\ii-c\jj-d\kk}{a^2+b^2+c^2+d^2},\]

    mentre la moltiplicazione non è commutativa.

\end{example}

\begin{definition}
    Un anello commutativo che è anche un corpo si dice \textbf{campo}.
\end{definition}

\begin{example}
    Alcuni campi, tra i più importanti, sono $\QQ$, $\RR$, $\CC$ e $\ZZ/p\ZZ$ con
    $p$ primo.
\end{example}

\begin{definition}
    Un elemento $a \neq 0$ appartenente a un anello $A$ si dice \textbf{divisore di zero} se
    $\exists b \neq 0 \in A \mid ab = 0$ o $ba = 0$.
\end{definition}

\begin{example}
    $2$ è un divisore di zero in $\ZZ/6\ZZ$, infatti $2 \cdot 3 \equiv 0 \pmod 6.$
\end{example}

\begin{definition}
    Un anello commutativo in cui non sono presenti divisori di zero si dice \textbf{dominio d'integrità},
    o più semplicemente \textit{dominio}.
\end{definition}

\begin{proposition}[\textit{Legge di annullamento del prodotto}]
    Sia $D$ un dominio. Allora $ab=0 \implies a=0 \,\lor\, b=0$.
\end{proposition}

\begin{proof}
    Siano $a$, $b \in D \mid ab = 0$. Se $a=0$, la condizione è soddisfatta.
    Se invece $a \neq 0$, $b$ deve essere per forza nullo, altrimenti si
    sarebbe trovato un divisore di $0$, e $D$ non sarebbe un dominio, \Lightning.
\end{proof}

\begin{example}
    L'anello dei polinomi su un campo, $\KK[x]$, è un dominio.
\end{example}

\section{Omomorfismi di anelli e ideali}

\begin{definition}
    Un \textbf{omomorfismo di anelli}\footnote{La specificazione "di anelli" è d'ora in avanti omessa.} è una mappa $\phi : A \to B$ -- con
    $A$ e $B$ anelli -- soddisfacente alcune particolari proprietà:

    \begin{itemize}
        \item $\phi$ è un \textit{omomorfismo di gruppi} rispetto all'addizione
              di $A$ e di $B$, ossia $\forall a, b \in A, \, \phi(a+b)=\phi(a)+\phi(b)$,
        \item $\phi(ab)=\phi(a)\phi(b)$,
        \item $\phi(1_A)=1_B$.
    \end{itemize}
\end{definition}

\begin{definition}
    Se $\phi : A \to B$ è un omomorfismo iniettivo, si dice che
    $\phi$ è un \textbf{monomorfismo}.
\end{definition}

\begin{definition}
    Se $\phi : A \to B$ è un omomorfismo suriettivo, si dice che
    $\phi$ è un \textbf{epimorfismo}.
\end{definition}

\begin{definition}
    Se $\phi : A \to B$ è un omomorfismo bigettivo\footnote{Ovvero se è sia un monomorfismo che un epimorfismo.}, si dice che
    $\phi$ è un \textbf{isomorfismo}.
\end{definition}

Prima di enunciare l'analogo del \textit{Primo teorema d'isomorfismo} dei gruppi
in relazione agli anelli, si rifletta su un esempio di omomorfismo:

\begin{example}
    Sia $\phi : \ZZ \to \ZZ, k \mapsto 2k$ un omomorfismo. Esso è un monomorfismo,
    infatti $\phi(x)=\phi(y) \implies 2x=2y \implies x=y$. Pertanto $\Ker \phi = \{0\}$. Sebbene $\Ker \phi < \ZZ$, esso \textbf{non è un sottoanello}\footnote{Infatti $1 \notin \Ker \phi$.}.
\end{example}

Dunque, con lo scopo di definire meglio le proprietà di un \textit{kernel},
così come si introdotto il concetto di \textit{sottogruppo normale} per i gruppi, si introduce ora il concetto di \textbf{ideale}.

\begin{definition}
    Si definisce ideale rispetto all'anello $A$ un insieme $I$ avente le seguenti proprietà:

    \begin{itemize}
        \item $I \leq A$,
        \item $\forall a \in A$, $\forall b \in I$, $ab \in I$ e $ba \in I$.
    \end{itemize}
\end{definition}

\begin{example}
    \label{exmpl:polinomi}
    Sia $I$ l'insieme dei polinomi di $\RR[x]$ tali che $2$ ne sia radice. Esso
    altro non è che un ideale, infatti $0 \in I \,\land\, \forall f(x), g(x) \in I, (f+g)(2)=0$ (i.e. $I<\RR[x]$) e $\forall f(x) \in A, \, g(x) \in I, \, (fg)(2) = 0$.
\end{example}

\begin{proposition}
    Sia $I$ un ideale di $A$. $1 \in I \implies I = A$.
\end{proposition}

\begin{proof}
    Per le proprietà dell'ideale $I$, $\forall a \in A$, $a1 = a \in I \implies
        A \subseteq I$. Dal momento che anche $I \subseteq A$, si deduce che $I = A$.
\end{proof}

\begin{proposition}
    Sia $\phi : A \to B$ un omomorfismo. $\Ker \phi$ è allora un ideale di $A$.
\end{proposition}

\begin{proof}
    Poiché $\phi$ è anche un omomorfismo tra gruppi, si deduce che $\Ker \phi \leq A$.
    Inoltre $\forall a \in A$, $\forall b \in \Ker \phi$, $\phi(ab)=\phi(a)\phi(b)=\phi(a)0=0 \implies ab \in I$.
\end{proof}

\begin{proposition}
    Sia $\phi : A \to B$ un omomorfismo. $\Imm \phi$ è allora un sottoanello di $B$.
\end{proposition}

\begin{proof}
    Chiaramente $0, 1 \in \Imm \phi$, dal momento che $\phi(0) = 0,\, \phi(1)=1$. Inoltre, dalla teoria dei gruppi, si ricorda anche che $\Imm \phi \leq B$.
    Infine, $\forall \phi(a),\, \phi(b) \in \Imm \phi, \, \phi(a)\phi(b) = \phi(ab) \in \Imm \phi$.
\end{proof}

\begin{definition}
    Si definisce con la notazione $(a)$ l'ideale \textit{bilatero} generato da $a$ in $A$, ossia:

    \[(a)=\{ba \mid b \in A\} \cup \{ab \mid b \in A\}.\]
\end{definition}

\begin{definition}
    Si dice che un ideale $I$ è \textit{principale} o \textbf{monogenerato}, quando $\exists a \in I \mid I = (a)$.
\end{definition}

\begin{example}
    In relazione all'\textit{Esempio \ref{exmpl:polinomi}}, l'ideale $I$ è
    monogenerato\footnote{Non è un caso: $\RR[x]$, in quanto anello euclideo, si dimostra essere un PID (\textit{principal ideal domain}), ossia un dominio che ammette \textit{solo} ideali monogenerati.}. In particolare, $I=(x-2)$.
\end{example}

\section{Quoziente per ideale e primo teorema d'isomorfismo}

Si definisce invece adesso il concetto di \textbf{anello quoziente}, in modo
completamente analogo a quello di \textit{gruppo quoziente}:

\begin{definition}
    Sia $A$ un anello e $I$ un suo ideale, si definisce $A/I$ l'anello ottenuto
    quozientando $A$ per $I$. Gli elementi di tale anello sono le classi di equivalenza di $\sim$ (i.e. gli elementi di $A/{\sim}$), dove $\forall a$, $b \in A$, $a\sim b \iff a-b \in I$. Tali classi di equivalenza vengono indicate come
    $a + I$, dove $a$ è un rappresentante della classe. L'anello è così dotato di due operazioni:

    \begin{itemize}
        \item $\forall a$, $b \in A$, $(a+I)+(b+I)=(a+b)+I$,
        \item $\forall a$, $b \in A$, $(a+I)(b+I)=ab+I$.
    \end{itemize}
\end{definition}

\begin{remark*}
    L'addizione di $A/I$ è ben definita, dal momento che $I \nsg A$, in quanto sottogruppo di un gruppo abeliano.
\end{remark*}

\begin{remark*}
    Anche la moltiplicazione di $A/I$ è ben definita. Siano $a\sim a'$, $b \sim b'$ quattro elementi di $A$ tali che $a = a' + i_1$ e $b = b' + i_2$ con $i_1$, $i_2 \in I$. Allora $ab=(a'+i_1)(b'+i_2)=a'b' + \underbrace{i_1b' + i_2a' + i_1i_2}_{\in I} \implies ab \sim a'b'$.
\end{remark*}

\begin{proposition}
    \label{prop:quoziente_pieno}
    $A/\{0\} \cong A$.
\end{proposition}

\begin{proof}
    Sia $\pi : A \to A/\{0\}$, $a \mapsto a + \{0\}$ l'omomorfismo di proiezione
    al quoziente. Innanzitutto, $a \sim a' \iff a-a'=0 \iff a=a'$, per cui $\pi$ è
    un monomorfismo (altrimenti si troverebbero due $a$, $b \mid a \neq b \,\land\, a \sim b$). Infine, $\pi$ è un epimorfismo, dal momento che $\forall a + \{0\} \in A/\{0\}, \, \pi(a) = a + \{0\}$. Pertanto $\pi$ è un isomorfismo.
\end{proof}

Adesso è possibile enunciare il seguente fondamentale teorema:

\begin{theorem}[\textit{Primo teorema d'isomorfismo}]
    \label{th:primo_isomorfismo}
    Sia $\phi : A \to B$ un omomorfismo. $A/\Ker \phi \cong \Imm \phi$.
\end{theorem}

\begin{proof}
    La dimostrazione procede in modo analogo a quanto visto per il teorema correlato
    in teoria dei gruppi. \\

    Sia $\zeta : A/\Ker \phi \to \Imm \phi$, $a + \Ker \phi \mapsto \phi(a)$.
    Si verifica che $\zeta$ è un omomorfismo: essendolo già per i
    gruppi, è sufficiente verificare che $\zeta((a+I)(b+I))=\zeta(ab+I)=\phi(ab)=\phi(a)\phi(b)=\zeta(a+I)\zeta(b+I)$. \\

    $\zeta$ è chiaramente anche un epimorfismo, dal momento che $\forall \phi(a) \in \Imm \phi$, $\zeta(a + \Ker \phi) = \phi(a)$. Inoltre, dal momento che $\zeta(a + \Ker \phi) = 0 \iff \phi(a) = 0 \iff a + \Ker \phi = \Ker \phi$, ossia l'identità di $A/\Ker \phi$, si deduce anche che $\zeta$ è un monomorfismo. Pertanto $\zeta$ è un isomorfismo.
\end{proof}

\begin{corollary}
    \label{cor:primo_isomorfismo_iniettivo}
    Sia $\phi : A \to B$ un monomorfismo. $A \cong \Imm \phi$.
\end{corollary}

\begin{proof}
    Poiché $\phi$ è un monomorfismo, $\Ker \phi = \{0\}$. Allora, per il \textit{Primo teorema di isomorfismo}, $A/\{0\} \cong \Imm \phi$. Dalla
    \textit{Proposizione \ref{prop:quoziente_pieno}}, si desume che $A \cong A/\{0\}$. Allora, per la proprietà transitiva degli isomorfismi, $A \cong \Imm \phi$.
\end{proof}

\end{document}

\newpage
\thispagestyle{empty}
~\newpage

\section{Anelli euclidei, PID e UFD}

\subsection{Prime proprietà}

Nel corso della storia della matematica, numerosi studiosi hanno tentato
di generalizzare -- o meglio, accomunare a più strutture algebriche -- il
concetto di divisione euclidea che era stato formulato per l'anello
dei numeri interi $\ZZ$ e, successivamente, per l'anello dei polinomi
$\KK[x]$. Lo sforzo di questi studiosi ad oggi è converso in un'unica
definizione, quella di anello euclideo, di seguito presentata.

\begin{definition}
    Un \textbf{anello euclideo} è un dominio d'integrità $D$\footnote{Difatti, nella
        letteratura inglese, si parla di \textit{Euclidean domain} piuttosto che di
        anello.} sul quale è
    definita una funzione $g$ detta \textbf{funzione grado} o \textit{norma}
    soddisfacente le seguenti proprietà:

    \begin{itemize}
        \item $g : D \setminus \{0\} \to \NN$,
        \item $\forall a$, $b \in D \setminus \{0\}$, $g(a) \leq g(ab)$,
        \item $\forall a \in D$, $b \in D \setminus \{0\}$, $\exists q$, $r \in D \mid
                  a=bq+r$ e $r=0 \,\lor\, g(r)<g(q)$.
    \end{itemize}
\end{definition}

Di seguito vengono presentate alcune definizioni, correlate
alle proprietà immediate di un anello euclideo.

\begin{definition}
    Dato un anello euclideo $E$, siano $a \in E$ e $b \in E \setminus \{0\}$. Si dice che
    $b \mid a$, ossia che $b$ \textit{divide} $a$, se $\exists c \in E \mid
        a=bc$.
\end{definition}

\begin{remark*}
    Si osserva che, per ogni anello euclideo $E$, qualsiasi $a \in E$ divide
    $0$. Infatti, $0 = a0$.
\end{remark*}

\begin{proposition}
    Dato un anello euclideo $E$, $a \mid b \,\land\, b \nmid a \implies g(a) < g(b)$.
\end{proposition}

\begin{proof}
    Poiché $b \nmid a$, esistono $q$, $r$ tali che $a = bq + r$, con
    $g(r) < g(b)$. Dal momento però che $a \mid b$, $\exists c \mid
        b = ac$. Pertanto $a = ac + r \implies r = a(1-c)$. Dacché $1-c \neq 0$ --
    altrimenti $r=0$, \Lightning{} --, così come $a \neq 0$, si deduce
    dalle proprietà della funzione grado che $g(a) \leq g(r)$.
    Combinando le due disuguaglianze, si ottiene la
    tesi: $g(a) < g(b)$.
\end{proof}

\begin{proposition}
    \label{prop:g1_minimo}
    $g(1)$ è il minimo di $\Imm g$, ossia il minimo grado assumibile
    da un elemento di un anello euclideo $E$.
\end{proposition}

\begin{proof}
    Sia $a \in E \setminus \{0\}$, allora, per le proprietà della funzione
    grado, $g(1) \leq g(1a) = g(a)$.
\end{proof}

\begin{theorem}
    Sia $a \in E \setminus \{0\}$, allora $a \in E^* \iff g(a) = g(1)$.
\end{theorem}

\begin{proof}
    Dividiamo la dimostrazione in due parti, ognuna corrispondente a una implicazione. \\

    ($\implies$) \;Sia $a \in E^*$, allora $\exists b \in E^*$ tale che $ab=1$. Poiché
    sia $a$ che $b$ sono diversi da $0$, dalle proprietà della funzione grado si
    desume che $g(a) \leq g(ab) = g(1)$. Poiché, dalla \textit{Proposizione \ref{prop:g1_minimo}},
    $g(1)$ è minimo, si conclude che $g(a) = g(1)$. \\

    ($\;\Longleftarrow\;$) \;Sia $a \in E \setminus \{0\}$ con $g(a) = g(1)$. Allora
    esistono $q$, $r$ tali che $1 = aq + r$. Vi sono due possibilità: che $r$ sia $0$, o
    che $g(r) < g(a)$. Tuttavia, poiché $g(a)=g(1)$, dalla \textit{Proposizione \ref{prop:g1_minimo}} si desume che $g(a)$ è minimo, e quindi che
    $r$ è nullo. Si conclude quindi che $aq = 1$, e dunque che $a \in E^*$.
\end{proof}

\subsection{Irriducibili e prime definizioni}

Come accade nell'aritmetica dei numeri interi, anche in un dominio è possibile definire
una nozione di \textit{primo}. In un dominio possono essere tuttavia definiti due tipi di "primi",
gli elementi \textit{irriducibili} e gli elementi \textit{primi}.

\begin{definition}
    In un dominio $A$, si dice che $a \in A \setminus A^*$ è \textbf{irriducibile} se
    $\exists b$, $c \mid a=bc \implies b \in A^*$ o $c \in A^*$.
\end{definition}

\begin{remark*}
    Dalla definizione si escludono gli invertibili di $A$ per permettere
    di definire meglio il concetto di fattorizzazione in seguito. Infatti,
    se li avessimo inclusi, avremmo che ogni dominio sarebbe a fattorizzazione
    non unica, dal momento che $a=bc$ potrebbe essere scritto anche come
    $a=1bc$.
\end{remark*}

\begin{definition}
    Si dice che due elementi non nulli $a$, $b$ appartenenti a un anello euclideo
    $E$ sono \textbf{associati} se $a \mid b$ e $b \mid a$.
\end{definition}

\begin{proposition}
    \label{prop:associati}
    $a$ e $b$ sono associati $\iff \exists c \in E^* \mid a=bc$ e $a$, $b$ entrambi non nulli.
\end{proposition}

\begin{proof} Si dimostrano le due implicazioni separatamente. \\

    ($\implies$) Se $a$ e $b$ sono associati, allora $\exists d$, $e$ tali che $a=bd$ e che $b=ae$. Combinando le due relazioni si ottiene:

    \[ a=aed \implies a(1-ed)=0.\]

    Poiché $a$ è diverso da zero, si ricava che $ed=1$, ossia
    che $d$, $e \in E^*$, e quindi la tesi. \\

    ($\;\Longleftarrow\;$) Se $a$ e $b$ sono entrambi non
    nulli e $\exists c \in E^* \mid a=bc$, $b$ chiaramente
    divide $a$. Inoltre, $a=bc \implies b=ac^{-1}$, e quindi
    anche $a$ divide $b$. Pertanto $a$ e $b$ sono associati.
\end{proof}

\begin{proposition}
    \label{prop:divisione_associati}
    Siano $a$ e $b$ due associati in $E$. Allora $a \mid c \implies b \mid c$.
\end{proposition}

\begin{proof}
    Poiché $a$ e $b$ sono associati, per la \textit{Proposizione \ref{prop:associati}}, $\exists d \in E^*$ tale che
    $a = db$. Dal momento che $a \mid c$, $\exists \alpha \in E$ tale che
    $c = \alpha a$, quindi:

    \[ c = \alpha a = \alpha d b,\]

    da cui la tesi.
\end{proof}

\begin{proposition}
    \label{prop:associati_generatori}
    Siano $a$ e $b$ due associati in $E$. Allora
    $(a)=(b)$.
\end{proposition}

\begin{proof}
    Poiché $a$ e $b$ sono associati, $\exists d \in E^*$ tale che $a = db$. Si dimostra l'uguaglianza dei due insiemi. \\

    Sia $\alpha = ak \in (a)$, allora $\alpha = dbk$ appartiene anche a $(b)$, quindi $(a) \subseteq (b)$. Sia
    invece $\beta = bk \in (b)$, allora $\beta = d^{-1}ak$
    appartiene anche a $(a)$, da cui $(b) \subseteq (a)$.
    Dalla doppia inclusione si verifica la tesi, $(a)=(b)$.
\end{proof}

\begin{definition}
    In un dominio $A$, si dice che $a \in A \setminus A^*$ è \textbf{primo} se
    $a \mid bc \implies a \mid b \,\lor\, a \mid c$.
\end{definition}

\begin{proposition}
    Se $a \in A$ è primo, allora $a$ è anche irriducibile.
\end{proposition}

\begin{proof}
    Si dimostra la tesi contronominalmente. Sia $a$ non irriducibile. Se
    $a \in A^*$, allora $a$ non può essere primo. Altrimenti $a=bc$ con
    $b$, $c \in A \setminus A^*$. \\

    Chiaramente $a \mid bc$, ossia sé stesso. Senza perdità di generalità, se $a \mid b$, dal momento che anche $b \mid a$,
    si dedurrebbe che $a$ e $b$ sono associati secondo la
    \textit{Proposizione \ref{prop:associati}}. Tuttavia questo
    implicherebbe che $c \in A^*$, \Lightning{}.
\end{proof}

\subsection{PID e MCD}

Come accade per $\ZZ$, in ogni anello euclideo è possibile definire il
concetto di \textit{massimo comun divisore}, sebbene con qualche accortezza
in più. Pertanto, ancor prima di definirlo, si enuncia la definizione di
PID e si dimostra un teorema fondamentale degli anelli euclidei, che
si ripresenterà in seguito come ingrediente fondamentale per la fondazione
del concetto di MCD.

\begin{definition}
    Si dice che un dominio è un \textit{principal ideal domain} (\textbf{PID})\footnote{Ossia un \textit{dominio
            a soli ideali principali}, quindi monogenerati, proprio come da definizione.} se ogni suo ideale è monogenerato.
\end{definition}

\begin{theorem}
    Sia $E$ un anello euclideo. Allora $E$ è un PID.
\end{theorem}

\begin{proof}
    Sia $I$ un ideale di $E$. Se $I = (0)$, allora $I$ è già monogenerato.
    Altrimenti si consideri l'insieme $g(I \setminus \{0\})$. Poiché
    $g(I \setminus \{0\}) \subseteq \NN$,
    esso ammette un minimo per il principio del buon ordinamento. \\

    Sia $m \in I$ un valore che assume tale minimo e sia $a \in I$.
    Poiché $E$ è euclideo, $\exists q$, $r \mid a = mq + r$ con $r=0$ o
    $g(r)<g(m)$. Tuttavia, poiché $r = a-mg \in I$ e $g(m)$ è minimo, necessariamente $r=0$ -- altrimenti $r$ sarebbe
    ancor più minimo di $m$, \Lightning{} --,
    quindi $m \mid a$, $\forall a \in I$. Quindi $I \subseteq (m)$. \\

    Dal momento che per le proprietà degli ideali $\forall a \in E$, $ma \in I$,
    si conclude che $(m) \subseteq I$. Quindi $I = (m)$.
\end{proof}

Adesso è possibile definire il concetto di massimo comun divisore, basandoci
sul fatto che ogni anello euclideo è un PID.

\begin{definition}
    Sia $D$ un dominio e siano $a$, $b \in D$. Si definisce
    \textit{massimo comun divisore} (\textbf{MCD}) di $a$ e $b$ un
    generatore dell'ideale $(a,b)$.
\end{definition}

\begin{remark*}
    Questa definizione di MCD è una buona definizione dal momento che sicuramente
    esiste un generatore dell'ideale $(a,b)$, dacché $D$ è un PID.
\end{remark*}

\begin{remark*}
    Non si parla di un unico massimo comun divisore, dal momento che
    potrebbero esservi più generatori dell'ideale $(a,b)$. Segue tuttavia che tutti questi generatori sono in realtà
    associati\footnote{Infatti ogni generatore divide ogni
        altro elemento di un ideale, e così i vari generatori si
        dividono tra di loro. Pertanto sono associati.}.
    Quando si scriverà
    $\MCD(a,b)$ s'intenderà quindi uno qualsiasi di questi associati.
\end{remark*}

\begin{theorem}[\textit{Identità di Bézout}]
    \label{th:bezout}
    Sia $d$ un MCD di $a$ e $b$. Allora
    $\exists \alpha$, $\beta$ tali che $d = \alpha a + \beta b$.
\end{theorem}

\begin{proof}
    Il teorema segue dalla definizione di MCD come generatore
    dell'ideale $(a,b)$. Infatti, poiché $d \in (a,b)$, esistono
    sicuramente, per definizione, $\alpha$ e $\beta$ tali che
    $d = \alpha a + \beta b$.
\end{proof}

\begin{proposition}
    \label{prop:mcd}
    Siano $a$, $b \in D$. Allora vale la seguente equivalenza:

    \[ d = \MCD(a, b) \iff \begin{cases} d \mid a \,\land\, d \mid b \\ \forall c \text{ t.c.\ } c \mid a \,\land\, c \mid b,\;c \mid d\end{cases}\]
\end{proposition}

\begin{proof} Si dimostrano le due implicazioni separatamente. \\

    ($\implies$) Poiché $d$ è generatore dell'ideale $(a, b)$, la prima proprietà segue banalmente. \\

    Inoltre, per l'\nameref{th:bezout}, $\exists \alpha$, $\beta$ tali che
    $d = \alpha a + \beta b$. Allora, se $c \mid a$ e $c \mid b$, sicuramente
    esistono $\gamma$ e $\delta$ tali che $a=\gamma c$ e $b=\delta c$. Pertanto
    si verifica la seconda proprietà, e quindi la tesi:

    \[ d = \alpha a + \beta b = \alpha \gamma c + \beta \delta c = c(\alpha\gamma+\beta\delta). \]

    \vskip 0.1in

    ($\;\Longleftarrow\;$) Sia $m = \MCD(a,b)$. Dal momento che $d$ divide
    sia $a$ che $b$, $d$ deve dividere, per l'implicazione scorsa, anche $m$.
    Per la seconda proprietà, $m$ divide $d$ a sua volta. Allora $d$ è un
    associato di $m$, e quindi, dalla \textit{Proposizione \ref{prop:associati_generatori}}, $(m)=(d)=(a,b)$, da cui $d = \MCD(a,b)$.
\end{proof}

\begin{proposition}
    \label{prop:divisione_gcd}
    Se $a \mid bc$ e $d = \MCD(a, b) \in D^*$, allora $a \mid c$.
\end{proposition}

\begin{proof}
    Per l'\nameref{th:bezout} $\exists \alpha$, $\beta$ tali che
    $\alpha a + \beta b = d$. Allora, poiché $a \mid bc$, $\exists
        \gamma$ tale che $bc=a\gamma$. Si verifica quindi la tesi:

    \[ \alpha a + \beta b = d \implies \alpha ac + \beta bc = dc \implies
        a d^{-1} (\alpha c + \beta \gamma) = c.\]
\end{proof}

\begin{lemma}
    \label{lem:primalità_mcd}
    Se $a$ è un irriducibile di un PID $D$, allora $\forall b \in D$,
    $(a,b)=D \,\lor\, (a,b)=(a)$, o equivalentemente $\MCD(a,b) \in D^*$ o
    $\MCD(a,b) = a$.
\end{lemma}

\begin{proof}
    Dacché $\MCD(a,b) \mid a$, le uniche opzioni, dal momento che $a$ è irriducibile,
    sono che $\MCD(a,b)$ sia un invertibile o che sia un associato
    di $a$ stesso.
\end{proof}

\begin{theorem}
    \label{th:irriducibili_primi}
    Se $a$ è un irriducibile di un PID $D$, allora $a$ è anche un primo.
\end{theorem}

\begin{proof}
    Siano $b$ e $c$ tali che $a \mid bc$. Per il \textit{Lemma \ref{lem:primalità_mcd}},
    $\MCD(a,b)$ può essere solo un associato di $a$ o essere un invertibile. Se è
    un associato di $a$, allora, per la \textit{Proposizione \ref{prop:divisione_associati}}, poiché $\MCD(a,b)$ divide $b$, anche $a$ divide $b$.
    Altrimenti $\MCD(a,b) \in D^*$, e quindi, per la \textit{Proposizione \ref{prop:divisione_gcd}}, $a \mid c$.
\end{proof}

\subsection{L'algoritmo di Euclide}

Per algoritmo di Euclide si intende un algoritmo che è in grado di
produrre in un numero finito di passi un MCD tra due elementi
$a$ e $b$ non entrambi nulli di un anello euclideo\footnote{Si richiede che l'anello sia
    euclideo e non soltanto che sia un PID, dal momento che l'algoritmo
    usufruisce delle proprietà della funzione grado.}. L'algoritmo
classico è di seguito presentato:

\newpage

\begin{algorithm}
    $e \gets \max(a,b)$\;
    $d \gets \min(a,b)$\;
    \BlankLine\BlankLine
    \While{$d>0$}
    {
        $m \gets d$\;
        $d \gets e \bmod d$\;
        $e \gets m$\;
    }
\end{algorithm}

dove $e$ è l'MCD ricercato e l'operazione $\mathrm{mod}$ restituisce un resto della
divisione euclidea\footnote{Ossia $a \bmod b$ restituisce un $r$ tale che $\exists q
        \mid a = bq+r$ con $r=0$ o $g(r)<g(q)$.}.

\begin{lemma}
    \label{lem:euclide_finito}
    L'algoritmo di Euclide termina sempre in un numero finito di passi.
\end{lemma}

\begin{proof}
    Se $d$ è pari a $0$, l'algoritmo termina immediatamente. \\

    Altrimenti si può costruire una sequenza $(g(d_i))_{i\geq1}$ dove $d_i$ è il valore di $d$ all'inizio
    di ogni $i$-esimo ciclo $\textbf{while}$. Ad ogni ciclo vi sono due casi: se $d_i$ si annulla dopo
    l'operazione di $\mathrm{mod}$, il ciclo si conclude al passo successivo, altrimenti,
    poiché $d_i$ è un resto di una divisione euclidea, segue che $g(d_i)<g(d_{i-1})$, dove
    si pone $d_{0}=\min(a, b)$. \\

    Per il principio della discesa infinita, $(g(d_i))_{i\geq1}$ non può essere
    una sequenza infinita, essendo strettamente decrescente. Quindi la sequenza è
    finita, e pertanto il ciclo $\textbf{while}$ s'interrompe dopo un numero finito
    di passi.
\end{proof}

\begin{lemma}
    \label{lem:generatori_euclide}
    Sia $r = a \bmod b$. Allora vale che $(a,b)=(b,r)$.
\end{lemma}

\begin{proof}
    Poiché $r = a \bmod b$, $\exists q$ tale che $a = qb + r$.
    Siano $k_1$ e $k_2$ tali che $(k_1)=(a,b)$ e $(k_2)=(b,r)$. Dal
    momento che $k_1$ divide sia $a$ che $b$, si ha che divide anche
    $r$. Siano $\alpha$, $\beta$ tali che $a = \alpha k_1$ e
    $b = \beta k_1$. Si verifica infatti che:

    \[ r = a - qb = \alpha k_1 - q \beta k_1 = k_1 (\alpha - q \beta). \]

    Poiché $k_1$ divide sia $b$ che $r$, per le proprietà del $\MCD$,
    $k_1$ divide anche $k_2$. Analogamente, $k_2$ divide $k_1$. Pertanto
    $k_1$ e $k_2$ sono associati, e dalla \textit{Proposizione \ref{prop:associati_generatori}} generano quindi lo stesso ideale, da
    cui la tesi.
\end{proof}

\begin{theorem}
    L'algoritmo di Euclide restituisce sempre correttamente un MCD tra due elementi $a$ e $b$ non entrambi nulli in un numero finito di passi.
\end{theorem}

\begin{proof}
    Per il \textit{Lemma \ref{lem:euclide_finito}}, l'algoritmo sicuramente termina.
    Se $d$ è pari a $0$, allora l'algoritmo termina restituendo $e$. Il valore è
    corretto, dal momento che, senza perdità di generalità, se $b$ è nullo, allora
    $\MCD(a, b)=a$: infatti $a$ divide sia sé stesso che $0$, e ogni divisore di $a$ è
    sempre un divisore di $0$. \\

    Se invece $d$ non è pari a $0$, si scelga il $d_n$ tale che $g(d_n)$ sia l'ultimo
    elemento della sequenza $(g(d_i))_{i\geq1}$ definita nel \textit{Lemma \ref{lem:euclide_finito}}. Per il \textit{Lemma \ref{lem:generatori_euclide}},
    si ha la seguente uguaglianza:

    \[ (e_0, d_0) = (d_0, d_1) = \cdots = (d_n, 0) = (d_n). \]

    \vskip 0.1in

    Poiché quindi $d_n$ è generatore di $(e_0, d_0)=(a,b)$, $d_n = \MCD(a,b)$.
\end{proof}

\subsection{UFD e fattorizzazione}

Si enuncia ora la definizione fondamentale di UFD, sulla
quale costruiremo un teorema fondamentale per gli anelli
euclidei.

\begin{definition}
    Si dice che un dominio $D$ è uno \textit{unique factorization domain} (\textbf{UFD})\footnote{Ossia
        un \textit{dominio a fattorizzazione unica}.} se ogni $a \in D$ non nullo e non invertibile può essere scritto
    in forma unica come prodotto di irriducibili, a meno di associati.
\end{definition}

\begin{lemma}
    \label{lem:fattorizzazione}
    Sia $E$ un anello euclideo. Allora ogni elemento $a \in E$ non nullo e
    non invertibile può essere scritto come prodotto di irriducibili.
\end{lemma}

\begin{proof}
    Si definisca $A$ nel seguente modo:

    \[A = \{g(a) \mid a \in E \setminus (E^* \cup \{0\}) \text{ non sia prodotto di irriducibili}\}.\]

    \vskip 0.1in

    Se $A \neq \emptyset$, allora, poiché $A \subseteq \NN$, per il principio
    del buon ordinamento, esiste un $m \in E$ tale che $g(m)$ sia minimo.
    Sicuramente $m$ non è irriducibile -- altrimenti $g(m) \notin A$, \Lightning{} --,
    quindi $m=ab$ con $a$, $b \in E \setminus E^*$. \\

    Poiché $a \mid m$, ma $m \nmid a$ -- altrimenti $a$ e $m$ sarebbero
    associati, e quindi $b$ sarebbe invertibile --, si deduce che $g(a) < g(m)$, e
    quindi che $g(a) \notin A$. Allora $a$ può scriversi come prodotto di irriducibili.
    Analogamente anche $b$ può scriversi come prodotto di irriducibili, e quindi
    $m$, che è il prodotto di $a$ e $b$, è prodotto di irriducibili, \Lightning{}. \\

    Quindi $A = \emptyset$, e ogni $a \in E$ non nullo e non invertibile è prodotto
    di irriducibili.
\end{proof}

\begin{theorem}
    \label{th:euclidei_ufd}
    Sia $E$ un anello euclideo. Allora $E$ è un UFD\footnote{In realtà questo teorema
        è un caso particolare di un teorema più generale: ogni PID è un UFD. Poiché
        la dimostrazione esula dalle intenzioni di queste dispense, si è preferito
        dimostrare il caso più familiare. Per la dimostrazione del teorema più generale si
        rimanda a \cite[pp.~124-126]{di2013algebra}.}.
\end{theorem}

\begin{proof}
    Innanzitutto, per il \textit{Lemma \ref{lem:fattorizzazione}}, ogni
    $a \in E$ non invertibile e non nullo ammette una fattorizzazione. \\

    Sia allora $a \in E$ non invertibile e non nullo. Affinché $E$ sia un UFD,
    deve verificarsi la seguente condizione: se
    $a=p_1p_2 \cdots p_r=q_1q_2 \cdots q_s \in E$, allora
    $r=s$ ed esiste una permutazione $\sigma \in S_r$ tale per cui
    $\sigma$ associ a ogni indice $i$ di un $p_i$ un indice $j$ di
    un $q_j$ in modo tale che $p_i$ e $q_j$ siano associati. \\

    Si procede per induzione. \\

    \,(\textit{passo base}) \,Se $r=1$, allora $a$ è irriducibile. Allora necessariamente
    $s=1$, altrimenti $a$ sarebbe prodotto di irriducibili, e quindi contemporaneamente
    anche non irriducibile. Inoltre esiste la permutazione banale $e \in S_1$ che
    associa $p_1$ a $q_1$. \\

    \,(\textit{passo induttivo}) \,Si assume che valga la tesi se $a$ è
    prodotto di $r-1$ irriducibili.
    Si consideri $p_1$: poiché $p_1$ divide $a$, $p_1$ divide anche
    $q_1q_2 \cdots q_s$. Dal momento che $E$, in quanto
    anello euclideo, è anche un dominio, dal \textit{Teorema \ref{th:irriducibili_primi}}, $p_1$ è anche primo,
    e quindi $p_1 \mid q_1$ o $p_1 \mid q_2 \cdots q_s$. \\

    Se $p_1 \nmid q_1$ si reitera il procedimento su $q_2 \cdots q_s$, trovando in
    un numero finito di passi un $q_j$ tale per cui $p_1 \mid q_j$. Allora si procede
    la dimostrazione scambiando $q_1$ e $q_j$. \\

    Poiché $q_1$ è irriducibile, $p_1$ e $q_1$ sono associati, ossia $q_1 = kp_1$ con
    $k \in E^*$. Allora $p_1 \cdots p_r = q_1 \cdots q_s = kp_1 \cdots q_s$, quindi,
    dal momento che $p_1 \neq 0$ ed $E$ è un dominio:

    \[p_1(p_2 \cdots p_r - kq_2 \cdots q_s)=0 \implies p_2 \cdots p_r = kq_2 \cdots q_s .\]

    Tuttavia il primo membro è un prodotto $r-1$ irriducibili, pertanto $r=s$ ed
    esiste un $\sigma \in S_{r-1}$ che associa ad ogni irriducibile $p_i$ un suo
    associato $q_i$. Allora si estende $\sigma$ a $S_r$ mappando $p_1$ a $q_1$,
    verificando la tesi.
\end{proof}

\subsection{Il teorema cinese del resto}

    Il noto \nameref{th:cinese} è un risultato più generale di quanto
    si sia visto nel contesto dell'aritmetica modulare. Difatti, esso è
    applicabile in forma estesa a tutti gli anelli euclidei, non solo ai
    numeri interi (che comunque rimangono un esempio classico di anello euclideo). \\
    
    \begin{lemma}
        \label{lem:pre_cinese}

        Sia $a$ un elemento riducibile di un anello euclideo $E$ e
        sia $a=bc$, dove $\MCD(b, c) \in E^*$. Allora vale
        il seguente isomorfismo:
        
        \[ A/(a) \cong A/(b) \times A/(c). \]
    \end{lemma}
    
    \begin{proof}
        Si consideri la funzione $\pi$ definita nel seguente
        modo:
        
        \[ \pi : A/(a) \to A/(b) \times A/(c),\,e + (a) \mapsto (e + (b), e + (c)). \]
        
        \vskip 0.1in
        
        Si verifica che $\pi$ è un omomorfismo. Infatti
        $\pi(1 + (a)) = (1 + (b), 1 + (c))$. \\
        
        Siano $e$, $k \in A$. Allora
        $\pi$ soddisfa la linearità:
        
        \begin{multline*}
            \pi\Bigl(\bigl(e + (a)\bigr) + \bigl(k + (a)\bigr)\Bigr) = \pi\bigl(e + k + (a)\bigr) =
            \bigl(e + k + (b), e + k + (c)\bigr) =
            \bigl(e + (b), e + (c)\bigr) + \\ \bigl(k + (b), k + (c)\bigr) = \pi\bigl(e + (a)\bigr) +
            \pi\bigl(k + (a)\bigr).
        \end{multline*}
        
        \vskip 0.1in
        
        e la moltiplicatività:
        
        \begin{multline*}
            \pi\Bigl(\bigl(e + (a)\bigr) \cdot \bigl(k + (a)\bigr)\Bigr) = \pi\bigl(ek + (a)\bigr) =
            \bigl(ek + (b), ek + (c)\bigr) =
            \bigl(e + (b), e + (c)\bigr) \cdot \\ \bigl(k + (b), k + (c)\bigr) = \pi\bigl(e + (a)\bigr) \cdot
            \pi\bigl(k + (a)\bigr).
        \end{multline*}
        
        \vskip 0.1in
        
        Si studia $\Ker \pi$ per dimostrare l'iniettività di $\pi$.
        Si pone dunque $\pi\bigl(e + (a)\bigr) = \bigl(0 + (b), 0 + (c)\bigr)$.
        Questa condizione è equivalente ad asserire che sia $b$ che $c$ dividano
        $e$. \\
        
        Sia allora $k \in E$ tale che $e=bk$. Dal momento che $c$ divide $e$, si
        $e$ divide $bk$. Allora, dacché per ipotesi $\MCD(a, b) \in E^*$, per la
        \propref{prop:divisione_gcd} $c$ divide $k$. Quindi esiste $j \in E$ tale che
        $k=cj$. In particolare, unendo le due condizioni si ottiene
        $e=bk=bcj=aj$. Pertanto $a$ divide $e$, da cui si deduce che $e + (a)$
        è equivalente a $0 + (a)$. Allora, poiché $\Ker \pi = (0)$, $\pi$ è un
        monomorfismo. \\
        
        Si studia invece adesso la surgettività di $\pi$. Siano $\alpha$,
        $\beta \in E$. Si pone dunque $\pi\bigl(e + (a)\bigr) =
        \bigl(\alpha + (b), \beta + (c)\bigr)$. Questa condizione è equivalente
        al seguente sistema:
        
        \[ \begin{cases} e = \alpha + bk, \\ e = \beta + cj, \end{cases} \quad \text{con } k, j \in E. \]
        
        \vskip 0.1in
        
        Unendo le due condizioni si ottiene la seguente equazione:
        
        \[ \alpha + bk = \beta + cj \iff cj - bk = \alpha - \beta.  \]
        
        \vskip 0.1in
        
        Si consideri ora $d = \MCD(b, c)$. Per l'\nameref{th:bezout} esistono
        $x$, $y$ tali che:
        
        \[ cx+by=d, \]
        
        \vskip 0.1in
        
        da cui si ricava che:
        
        \[ (\alpha-\beta)(cx + by) = (\alpha-\beta)d \implies cxd\inv(\alpha-\beta)+byd\inv(\alpha-\beta)=\alpha-\beta, \]
        
        \vskip 0.1in
        
        ponendo allora $j=xd\inv(\alpha-\beta)$ e $k=-yd\inv(\alpha-\beta)$
        si ricava una possibile soluzione per $e$. Quindi
        $\pi$ è un epimorfismo. \\
        
        Poiché $\pi$ è sia un monomorfismo che un epimorfismo, si conclude
        che $\pi$ è un isomorfismo, da cui la tesi.
        
    \end{proof}

    \begin{theorem}[\textit{Teorema cinese del resto}]
        \label{th:cinese}
    
        Sia $a$ un elemento di un anello euclideo $A$ e sia
        $p_1^{m_1}p_2^{m_2}\cdots p_n^{m_n}$ una sua fattorizzazione
        in irriducibili non associati.
        Allora vale il seguente isomorfismo:
        
        \[ A/(a) \cong A/(p_1^{m_1}) \times \cdots \times A/(p_n^{m_n}). \]
    \end{theorem}
    
    \begin{proof}
        Si dimostra il teorema applicando il principio di induzione su $n$,
        il numero di fattori irriducibili distinti che appaiono
        nella fattorizzazione di $a$. \\
        
        \,(\textit{passo base}) \, Se $a$ consta di un solo fattore irriducibile,
        allora banalmente $A/(a) \cong A/(p_1^{m_1})$. \\
        
        \,(\textit{passo induttivo}) \, Possiamo riscrivere $a$ come
        il prodotto di $(p_1^{m_1}\cdots p_{n-1}^{m_{n-1}})$ e di
        $p_n^{m_n}$. \\
        
        Si nota innanzitutto che $d = \MCD(p_1^{m_1}\cdots p_{n-1}^{m_{n-1}}, p_n^{m_n})$
        è un invertibile. Se così non fosse, infatti, si potrebbe
        considerare un irriducibile $q$ della fattorizzazione di $d$:
        tale $q$, in quanto primo per il \thref{th:irriducibili_primi},
        deve dividere un $p_j$ con $1 \leq j \leq n-1$, così
        come deve dividere $p_n$. Allora $p_j$ e $q$ sono associati,
        così come $q$ e $p_n$. Dunque anche $p_j$ e $p_n$ sono associati.
        Tuttavia questo è un assurdo, dal momento che per ipotesi
        la fattorizzazione di $a$ include irriducibili distinti e
        non associati, \Lightning{}. \\
        
        Allora dal \lemref{lem:pre_cinese} si ricava che:
        
        \[ A/(a) \cong A/(p_1^{m_1}\cdots p_{n-1}^{m_{n-1}}) \times A/(p_n^{m_n}), \]
        
        \vskip 0.1in
        
        mentre dal passo induttivo si sa già che:
        
        \[ A/(p_1^{m_1}\cdots p_{n-1}^{m_{n-1}}) \cong A/(p_1^{m_1}) \times \cdots \times A/(p_{n-1}^{m_{n-1}}). \]
        
        \vskip 0.1in
        
        Pertanto, unendo le due informazioni, si verifica la tesi:
        
            \[ A/(a) \cong
            A/(p_1^{m_1}) \times \cdots \times A/(p_{n-1}^{m_{n-1}}) \times A/(p_n^{m_n}). \]
        
    \end{proof}


\newpage
\thispagestyle{empty}
~\newpage

% !BIB TS-program = biber

\PassOptionsToPackage{main=italian}{babel}
\documentclass[11pt]{scrbook}
\usepackage{evan_notes}
\usepackage[utf8]{inputenc}
\usepackage[italian]{babel}
\usepackage{algorithm2e}
\usepackage{amsfonts}
\usepackage{amsthm}
\usepackage{amssymb}
\usepackage{amsopn}
\usepackage[backend=biber]{biblatex}
\usepackage{cancel}
\usepackage{csquotes}
\usepackage{mathtools}
\usepackage{marvosym}

\addbibresource{bibliography.bib}

\begin{document}

\chapter{Esempi notevoli di anelli euclidei}

\section{I numeri interi: \texorpdfstring{$\ZZ$}{Z}}

Senza ombra di dubbio l'esempio più importante di anello euclideo -- nonché
l'esempio da cui si è generalizzata proprio la stessa nozione di anello
euclideo -- è l'anello dei numeri interi. \\

In questo dominio la funzione grado è canonicamente il valore assoluto:

\[g : \ZZ \setminus \{0\} \to \NN, \, k \mapsto \left|k\right|.\]

\vskip 0.1in

Infatti, chiaramente $|a| \leq |ab|\, \forall a$, $b \in \ZZ \setminus \{0\}$. Inoltre
esistono -- e sono anche unici, a meno di segno -- $q$, $r \in \ZZ \mid a = bq + r$, con $r=0 \,\lor\,
    \left|r\right| < \left|q\right|$. \\

Dal momento che così si verifica che $\ZZ$ è un anello euclideo, il \textit{Teorema
    fondamentale dell'aritmetica} è un corollario del teorema per cui ogni anello
    euclideo è un UFD.

\section{I campi: \texorpdfstring{$\KK$}{K}}

Ogni campo $\KK$ è un anello euclideo, seppur banalmente. Infatti, eccetto proprio
per $0$, ogni elemento è "divisibile" per ogni altro elemento: siano $a$, $b \in \KK$,
allora $a = ab^{-1}b$. \\

Si definisce quindi la funzione grado come la funzione nulla:

\[g : \KK^* \to \NN, \, a \mapsto 0.\]

\vskip 0.1in

Chiaramente $g$ soddisfa il primo assioma della funzione grado. Inoltre,
poiché ogni elemento è "divisibile", il resto è sempre zero -- non è pertanto
necessario verificare nessun'altra proprietà.

\section{I polinomi di un campo: \texorpdfstring{$\KK[x]$}{K[x]}}

I polinomi di un campo $\KK$ formano un anello euclideo rilevante
nello studio dell'algebra astratta. Come suggerisce la
terminologia, la funzione grado in questo dominio coincide
proprio con il grado del polinomio, ossia si definisce come:

\[g : \KK[x] \setminus \{0\} \to \NN, \, f(x) \mapsto \deg f.\]

\vskip 0.1in

Si verifica facilmente che $g(a(x)) \leq g(a(x)b(x)) \, \forall a(x)$, $b(x) \in \KK[x] \setminus \{0\}$, mentre la divisione euclidea -- come negli interi -- ci permette
di concludere che effettivamente $\KK[x]$ soddisfa tutti gli assiomi di un anello
euclideo\footnote{Curiosamente i polinomi di $\KK[x]$ e i campi $\KK$ sono gli unici anelli euclidei in cui resti
    e quozienti sono unici, includendo la scelta di segno (vd.
    \cite{10.2307/2315810}).}.

\begin{example}
    Sia $\alpha \in \KK$ e sia $\varphi_\alpha : \KK[x] \to \KK, \, f(x) \mapsto f(\alpha)$
    la sua valutazione polinomiale in $\KK[x]$. $\varphi_\alpha$ è un omomorfismo, il cui
    nucleo è rappresentato dai polinomi in $\KK[x]$ che hanno $\alpha$ come radice. Poiché
    $\KK[x]$ è un PID, $\Ker \varphi$ deve essere monogenerato. $x-\alpha \in \Ker \varphi$
    è irriducibile, e quindi è il generatore dell'ideale. Si deduce così che
    $\Ker \varphi = (x-\alpha)$.
\end{example}

\section{Gli interi di Gauss: \texorpdfstring{$\ZZ[i]$}{Z[i]}}

Un importante esempio di anello euclideo è il dominio degli interi di Gauss $\ZZ[i]$, definito come:

\[\ZZ[i] = \{a+bi \mid a, b \in \ZZ\}.\]

\vskip 0.1in

La funzione grado coincide in particolare con il quadrato del modulo di un numero complesso, ossia:
\[g(z) : \ZZ[i] \setminus \{0\} \to \NN, \, a+bi \mapsto \left| a+bi \right|^2.\]

Il vantaggio di quest'ultima definizione è l'enfasi sul collegamento tra la funzione grado
di $\ZZ$ e quella di $\ZZ[i].$ Infatti, se $a \in \ZZ$, il grado di $a$ in $\ZZ$ e in $\ZZ[i]$
sono uno il quadrato dell'altro. In particolare, è possibile ridefinire il grado
di $\ZZ$ proprio in modo tale da farlo coincidere con quello di $\ZZ[i]$. \\

\begin{theorem}
    $\ZZ[i]$ è un anello euclideo.
\end{theorem}

\begin{proof}
    Si verifica la prima proprietà della funzione grado. Siano $a$, $b \in \ZZ[i] \setminus \{0\}$,
    allora $\left|a\right| \geq 1 \,\land\, \left|b\right| \geq 1$. Poiché
    $\left|ab\right| = \left|a\right|\left|b\right|$\footnote{Questa interessante proprietà del modulo è alla base dell'identità di Brahmagupta-Fibonacci: $(a^2 + b^2)(c^2 + d^2) = (ac-bd)^2 + (ad+bc)^2.$}, si verifica facilmente che
    $\left|ab\right| \geq \left|a\right|$, ossia che $g(ab) \geq g(a)$. \\

    Si verifica infine che esiste una divisione euclidea, ossia che
    $\forall a \in \ZZ[i]$, $\forall b \in \ZZ[i] \setminus \{0\}$, $\exists q$, $r \in \ZZ[i] \mid a = bq + r$ e $r=0 \,\lor\, g(r) < g(b)$.
    Tutti i multipli di $b$ formano un piano con basi $b$ e $ib$, dove
    sicuramente esiste un certo $q$ tale che la distanza $\left|r\right| = \left|a-bq\right|$ sia minima. \\

    Se $a$ è un multiplo di $b$, vale sicuramente che $a = bq$. Altrimenti dal momento che $r$ è sicuramente inquadrato in uno dei tasselli del piano, vale
    sicuramente la seguente disuguaglianza, che lega il modulo di $r$ alla diagonale di
    ogni quadrato:

    \[\left|r\right| \leq \frac{\left|b\right|}{\sqrt{2}}.\]

    Pertanto vale la seconda e ultima proprietà della funzione grado:

    \[\left|r\right| \leq \frac{\left|b\right|}{\sqrt{2}} < \left|b\right| \implies \left|r\right|^2 < \left|b\right|^2 \implies g(r) < g(b).\]
\end{proof}

\section{Gli interi di Eisenstein: \texorpdfstring{$\ZZ[\omega]$}{Z[ω]}}

Sulla scia di $\ZZ[i]$ è possibile definire anche l'anello degli
interi di Eisenstein, aggiungendo a $\ZZ$ la prima radice cubica
primitiva dell'unità in senso antiorario, ossia:

\[\omega = e^{\frac{2\pi i}{3}} = -\frac{1}{2} + \frac{\sqrt{3}}{2}i.\]

In particolare, $\omega$ è una delle due radici dell'equazione
$z^2 + z + 1 = 0$, dove invece l'altra radice altro non è che
$\omega^2 = \overline{\omega}$.

\vskip 0.1in

La funzione grado in $\ZZ[\omega]$ deriva da quella di $\ZZ[i]$ e coincide ancora
con il quadrato del modulo del numero complesso. Si definisce quindi:

\[g : \ZZ[\omega] \setminus \{0\}, \, a+b\omega \mapsto \left|a+b\omega\right|^2.\]

Sviluppando il modulo è possibile ottenere una formula più concreta:

\[ \left|a+b\omega\right|^2 = \left|\left(a-\frac{b}{2}\right) + \frac{b\sqrt{3}}{2}i\right|^2 =\] \\

\[= \left(a-\frac{b}{2}\right)^2 + \frac{3b^2}{4} = a^2 - ab + b^2.\] \\

\begin{theorem}
    $\ZZ[\omega]$ è un anello euclideo.
\end{theorem}

\begin{proof}
    Sulla scia della dimostrazione presentata per $\ZZ[i]$, si verifica facilmente
    la prima proprietà della funzione grado. Siano $a$, $b \in \ZZ[\omega]$, allora
    $\left|a\right| \geq 1$ e $\left|b\right| \geq 1$. Poiché dalle proprietà
    dei numeri complessi vale ancora $\left|a\right| \left|b\right| \geq \left|a\right|$,
    la proprietà $g(ab) \geq g(a)$ è già verificata. \\

    Si verifica infine la seconda e ultima proprietà della funzione grado. Come per
    $\ZZ[i]$, i multipli di $b \in \ZZ[\omega]$ sono visualizzati su un piano che
    ha per basi $b$ e $\omega b$, pertanto
    esiste sicuramente un $q$ tale che la distanza $\left|a-bq\right|$ sia minima. \\

    Se $a$ è multiplo di $b$, allora chiaramente $a = bq$. Altrimenti, $a$ è certamente
    inquadrato in uno dei triangoli del piano, per cui vale la seguente disuguaglianza:

    \[\left|r\right| \leq \frac{\sqrt{3}}{2} \left|b\right|.\]

    Dunque la tesi è verificata:

    \[\left|r\right| \leq \frac{\sqrt{3}}{2} \left|b\right| < \left|b\right| \implies \left|r\right|^2 < \left|b\right|^2 \implies g(r) < g(b). \]
\end{proof}

\chapter{Riferimenti bibliografici}
\printbibliography[heading=none]

\end{document}

\newpage
\thispagestyle{empty}
~\newpage

% !BIB TS-program = biber

\PassOptionsToPackage{main=italian}{babel}
\documentclass[11pt]{scrbook}
\usepackage{evan_notes}
\usepackage[utf8]{inputenc}
\usepackage[italian]{babel}
\usepackage{algorithm2e}
\usepackage{amsfonts}
\usepackage{amsthm}
\usepackage{amssymb}
\usepackage{amsopn}
\usepackage[backend=biber]{biblatex}
\usepackage{cancel}
\usepackage{csquotes}
\usepackage{mathtools}
\usepackage{marvosym}

\addbibresource{bibliography.bib}

\begin{document}

\chapter{Irriducibili e corollari di aritmetica in \texorpdfstring{$\ZZi$}{Z[i]}}

Come già dimostrato, $\ZZi$ è un anello euclideo con la seguente
funzione grado:

\[ g : \ZZi \setminuszero \to \ZZ,\, a+bi \mapsto  \norm{a+bi}^2.\]

A partire da questo preconcetto è possibile dimostrare un teorema
importante in aritmetica, il \nameref{th:teorema_natale},
che discende direttamente come corollario di un teorema più
generale riguardante $\ZZi$.

\section{Il teorema di Natale di Fermat e gli irriducibili in \texorpdfstring{$\ZZi$}{Z[i]}}

\begin{lemma}
    \label{lem:riducibile_due_quadrati}
    Sia $p$ un numero primo riducibile in $\ZZi$, allora $p$
    può essere scritto come somma di due quadrati in $\ZZ$.
\end{lemma}

\begin{proof}
    Se $p$ è riducibile in $\ZZi$, allora esistono $a+bi$ e
    $c+di$ appartenenti a $\ZZi \setminus \ZZi^*$ tali che $p=(a+bi)(c+di)$. \\

    Impiegando le proprietà dell'operazione di coniugio si
    ottiene la seguente equazione:

    \[ \overline{p}=p=(a-bi)(c-di) \implies p^2=p \overline{p} = (a^2+b^2)(c^2+d^2). \]

    Dal momento che $a+bi$ e $c+di$ non sono invertibili,
    i valori della funzione grado calcolati in essi sono strettamente
    maggiori del valore assunto nell'unità, ovverosia:

    \[ a^2+b^2>1, \qquad c^2+d^2>1. \]

    Allora devono per forza valere le seguenti equazioni:

    \[ p=a^2+b^2, \qquad p=c^2+d^2, \]

    da cui la tesi.
\end{proof}

\begin{lemma}
    \label{lem:quadrato_mod_4}
    Sia $p$ un numero primo tale che $p \equiv 1 \pmod4$. Allora
    esiste un $x \in \ZZ$ tale che $p \mid x^2+1$.
\end{lemma}

\begin{proof}
    Per il \textit{Teorema di Wilson}, $(p-1)! \equiv -1 \pmod p$.
    Attraverso varie manipolazioni algebriche si ottiene:

    \[-1 \equiv 1 \cdots \frac{p-1}{2} \cdot \frac{p+1}{2} \cdots (p-1) \equiv 1 \cdots \frac{p-1}{2} \left(-\frac{p-1}{2}\right) \cdots (-1) \equiv\]

    \[ \equiv (-1)^{\frac{p-1}{2}} \left(\left( \frac{p-1}{2} \right)!\right)^2 \equiv
        \left(\left( \frac{p-1}{2} \right)!\right)^2 \pmod p,
    \]

    \vskip 0.1in

    da cui con $x = \left( \frac{p-1}{2} \right)!$ si verifica la
    tesi.
\end{proof}

\begin{theorem}
    \label{th:primo_1_mod_4_riducibile}
    Sia $p$ un numero primo tale che $p \equiv 1 \pmod4$. Allora
    $p$ è riducibile in $\ZZi$.
\end{theorem}

\begin{proof}
    Per il \textit{Lemma \ref{lem:quadrato_mod_4}}, si ha che esiste
    un $x \in \ZZ$ tale che $p \mid x^2+1$. Se $p$ fosse irriducibile,
    dacché $\ZZi$ è un PID in quanto euclideo, $p$ sarebbe anche un
    primo di $\ZZi$. Dal momento che $x^2+1=(x+i)(x-i)$, $p$ dovrebbe
    dividere almeno uno di questi due fattori. \\

    Senza perdità di generalità, si ponga che $p \mid (x+i)$. Allora
    $\exists a+bi \in \ZZi \mid x+i=(a+bi)p$. Uguagliando le parti
    immaginarie si ottiene $bp=1$, che non ammette soluzioni, \Lightning{}. Pertanto $p$ è riducibile.
\end{proof}

\begin{corollary}[\textit{Teorema di Natale di Fermat}]
    \label{th:teorema_natale}
    Sia $p$ un numero primo tale che $p \equiv 1 \pmod4$. Allora
    $p$ è somma di due quadrati in $\ZZ$.
\end{corollary}

\begin{proof}
    Per il \textit{Teorema \ref{th:primo_1_mod_4_riducibile}},
    $p$ è riducibile in $\ZZi$. In quanto riducibile in $\ZZi$, per
    il \textit{Lemma \ref{lem:riducibile_due_quadrati}}, $p$ è allora
    somma di due quadrati.
\end{proof}

\begin{theorem}
    \label{th:primo_-1_mod_4_irriducibile}
    Sia $p$ un numero primo tale che $p \equiv -1 \pmod4$. Allora
    $p$ è irriducibile in $\ZZi$.
\end{theorem}

\begin{proof}
    Se $p$ fosse riducibile in
    $\ZZi$, per il \nameref{th:teorema_natale} esisterebbero $a$ e $b$
    in $\ZZ$ tali che $p=a^2+b^2$. Dal momento che $p$ è dispari,
    possiamo supporre, senza perdità di generalità, che
    $a$ sia pari e che $b$ sia dispari. Pertanto $a^2 \equiv 0 \pmod 4$ e $b^2 \equiv 1 \pmod 4$, dacché sono uno pari e l'altro dispari\footnote{Infatti, $0^2 \equiv 0
            \pmod4$, $1^2 \equiv 1 \pmod4$, $2^2 \equiv 4 \equiv 0 \pmod 4$,
        $3^2 \equiv 9 \equiv 1 \pmod 4$.}. Tuttavia la congruenza
    $a^2+b^2 \equiv 1 \equiv -1 \pmod4$ non è mai soddisfatta,
    \Lightning{}. Pertanto $p$ può essere solo irriducibile.
\end{proof}

\begin{remark*}
    Si osserva che $2=(1+i)(1-i)$. Dal momento che $\norm{1+i}^2=
        \norm{1-i}^2=2\neq1$, si deduce che nessuno dei due fattori
    è invertibile. Pertanto $2$ non è irriducibile.
\end{remark*}

\begin{proposition}
    \label{prop:irriducibili_zz_zzi}
    Gli unici primi $p \in \ZZ$ irriducibili in $\ZZi$ sono i primi $p$ tali
    che $p \equiv -1 \pmod4$.
\end{proposition}

\begin{proof}
    Per l'osservazione precedente, $2$ non è irriducibile in $\ZZi$,
    così come i primi congrui a $1$ in modulo $4$,
    per il \textit{Teorema \ref{th:primo_1_mod_4_riducibile}}. Al
    contrario i primi $p$ congrui a $-1$ in modulo $4$ sono
    irriducibili, per il \textit{Teorema \ref{th:primo_-1_mod_4_irriducibile}}, da cui la tesi.
\end{proof}

\begin{theorem}
    $z \in \ZZi$ è irriducibile se e solo se $z$ è un associato di un primo $p \in \ZZ$ tale che $p \equiv -1 \pmod 4$, o se $\norm{z}^2$ è primo.
\end{theorem}

\begin{proof} Si dimostrano le due implicazioni separatamente. \\

    ($\implies$)\; Sia $z \in \ZZi$ irriducibile. Chiaramente
    $z \mid z \overline{z} = g(z)$. Dacché $\ZZ$ è un UFD,
    $g(z)$ può decomporsi in un prodotto di primi $q_1q_2\cdots q_n$.
    Dal momento che $\ZZi$ è un PID, in quanto anello euclideo,
    $z$ deve dividere uno dei primi della fattorizzazione di
    $g(z)$. Si assuma che tale primo sia $q_i$. Allora esiste
    un $w \in \ZZi$ tale che $q_i=wz$. \\

    Se $w \in \ZZi^*$, si
    deduce che $z$ è un associato di $q_i$. Dal momento che
    $z$ è irriducibile, $q_i$, che è suo associato, è a sua
    volta irriducibile. Allora, per la \textit{Proposizione \ref{prop:irriducibili_zz_zzi}}, $q_i \equiv -1 \pmod4$.
    \\

    Altrimenti, se $w$ non è invertibile, si ha che $g(w)>g(1)$,
    ossia che $\norm{w}^2>1$. Inoltre in quanto irriducibile, anche
    $z$ non è invertibile, e quindi
    $g(z)>g(1) \implies \norm{z}^2>1$. Dalla proprietà
    moltiplicativa
    del modulo si ricava $q_i^2 = \norm{q_i}^2 = \norm{w}^2 \norm{z}^2$,
    da cui necessariamente consegue che:

    \[ \norm{w}^2=q_i, \quad \norm{z}^2=q_i, \]

    attraverso cui si verifica l'implicazione. \\

    ($\,\Longleftarrow\,\,$)\; Se $p \in \ZZ$ e $p \equiv -1 \pmod4$, per
    il \textit{Teorema \ref{th:primo_-1_mod_4_irriducibile}}, $p$ è
    irriducibile. Allora in quanto suo associato, anche $z$ è irriducibile. \\

    Altrimenti, se $\norm{z}^2$ è un primo $p$, si ponga
    $z=ab$ con $a$ e $b \in \ZZi$. Per la proprietà moltiplicativa
    del modulo, $p = \norm{z}^2 = \norm{ab}^2 = \norm{a}^2\norm{b}^2$.
    Tuttavia questo implica che uno tra $\norm{a}^2$ e $\norm{b}^2$
    sia pari a $1$, ossia che uno tra $a$ e $b$ sia invertibile,
    dacché $g(1)=1$. Pertanto $z$ è in ogni caso irriducibile.
\end{proof}

Infine si enuncia un'ultima identità inerente all'aritmetica, ma
strettamente collegata a $\ZZi$.

\section{L'identità di Brahmagupta-Fibonacci}

\begin{proposition}[\textit{Identità di Brahmagupta-Fibonacci}]
    \label{prop:fibonacci}
    Il prodotto di due somme di quadrati è ancora una
    somma di quadrati. In particolare:

    \[ (a^2+b^2)(c^2+d^2)=(ac-bd)^2+(ad+bc)^2. \]
\end{proposition}

\begin{proof}
    La dimostrazione altro non è che una banale verifica
    algebrica. Ciononostante è possibile risalire a questa
    identità in via alternativa mediante l'uso
    del modulo dei numeri complessi. \\

    Siano $z_1=a+bi$, $z_2=c+di \in \CC$. Allora, per le proprietà
    del modulo dei numeri complessi:

    \begin{equation}
        \label{eq:modulo_z}
        \norm{z_1}\norm{z_2}=\norm{z_1z_2}.
    \end{equation}


    Computando il prodotto tra $z_1$ e $z_2$ si ottiene:

    \[ z_1z_2 = (ac-bd) + (ad+bc)i, \]

    da cui a sua volta si ricava:

    \[ \norm{z_1z_2} = \sqrt{(ac-bd)^2 + (ad+bc)^2}, \]

    assieme a:

    \[ \norm{z_1}=\sqrt{a^2+b^2}, \quad \norm{z_2}=\sqrt{c^2+d^2}. \]

    Infine, da \eqref{eq:modulo_z}, elevando al quadrato, si deduce l'identità
    presentata:

    \begin{multline*}
        \sqrt{a^2+b^2}\sqrt{c^2+d^2}=\sqrt{(ac-bd)^2 + (ad+bc)^2} \implies (a^2+b^2)(c^2+d^2)= \\ (ac-bd)^2+(ad+bc)^2.
    \end{multline*}
\end{proof}

\begin{example}
    Si consideri $65=5 \cdot 13$. Dal momento che sia $5$
    che $13$ sono congrui a $1$ in modulo $4$, sappiamo
    già si possono scrivere entrambi come somme di due
    quadrati. Allora, dall'\nameref{prop:fibonacci},
    anche $65$ è somma di due quadrati. \\

    Infatti $5=2^2+1^2$ e $13=3^2+2^2$. Pertanto
    $65=5\cdot 13=(2\cdot3-1\cdot2)^2 + (2\cdot2+1\cdot3)^2=4^2+7^2$.
\end{example}

\end{document}


\newpage
\thispagestyle{empty}
~\newpage

% !BIB TS-program = biber

\PassOptionsToPackage{main=italian}{babel}
\documentclass[11pt]{scrbook}
\usepackage{evan_notes}
\usepackage[utf8]{inputenc}
\usepackage[italian]{babel}
\usepackage{algorithm2e}
\usepackage{amsfonts}
\usepackage{amsthm}
\usepackage{amssymb}
\usepackage{amsopn}
\usepackage[backend=biber]{biblatex}
\usepackage{cancel}
\usepackage{csquotes}
\usepackage{mathtools}
\usepackage{marvosym}

\begin{document}

\chapter{Irriducibilità in \texorpdfstring{$\ZZx$}{Z[x]} e in \texorpdfstring{$\QQx$}{Q[x]}}

\section{Criterio di Eisenstein e proiezione in \texorpdfstring{$\ZZpx$}{Z\_p[x]}}

Prima di studiare le irriducibilità in $\ZZ$, si guarda
alle irriducibilità nei vari campi finiti $\ZZp$, con
$p$ primo. Questo metodo presenta un vantaggio da non
sottovalutare: in $\ZZp$ per ogni grado $n$ esiste un
numero finito di polinomi monici\footnote{Si prendono in
    considerazione solo i polinomi monici dal momento che vale
    l'equivalenza degli associati: se $a$ divide $b$, allora
    tutti gli associati di $a$ dividono $b$. $\ZZp$ è infatti
    un campo, e quindi $\ZZpx$ è un anello euclideo.} -- in particolare, $p^n$ --
e quindi per un polinomio di grado $d$ è sufficiente controllare
che questo non sia prodotto di tali polinomi monici per
$1 \leq n < d$. \\

In modo preliminare, si definisce un omomorfismo fondamentale.

\begin{definition}
    Sia il seguente l'\textbf{omomorfismo di proiezione} da
    $\ZZ$ in $\ZZp$:

    \[ \hatpip : \ZZx \to \ZZpx,\, a_n x^n + \ldots + a_0 \mapsto [a_n]_p \, x^n + \ldots + [a_0]_p. \]
\end{definition}

\begin{remark*}
    Si dimostra facilmente che $\hatpi$ è un omomorfismo di anelli.
    Innanzitutto, $\hatpi(1) = [1]_p$. Vale chiaramente la linearità:

    \begin{multline*}
        \hatpip(a_n x^n + \ldots + a_0) + \hatpip(b_n x^n + \ldots + b_0) = [a_n]_p \, x^n + \ldots + [b_n]_p \, x^n + \ldots = \\
        = [a_n+b_n]_p \, x^n + \ldots = \hatpip(a_n x^n + \ldots + a_0 + b_n x^n + \ldots + b_0).
    \end{multline*}

    Infine vale anche la moltiplicatività:

    \begin{multline*}
        \hatpip(a_n x^n + \ldots + a_0) \hatpip(b_n x^n + \ldots + b_0) = ([a_n]_p \, x^n + \ldots)([b_n]_p \, x^n + \ldots) = \\
        = \sum_{i=0}^n \sum_{j+k=i} [a_j]_p \, [b_k]_p \, x^i
        = \sum_{i=0}^n \sum_{j+k=i} [a_j b_k]_p \, x^i
        = \hatpip\left(\sum_{i=0}^n \sum_{j+k=i} a_j b_k x^i\right) = \\
        =\hatpip\left((a_n x^n + \ldots + a_0)(b_n x^n + \ldots + b_0)\right).
    \end{multline*}
\end{remark*}

Prima di enunciare un teorema che si rivelerà
importante nel determinare l'irriducibilità di un
polinomio in $\ZZx$, si enuncia una definizione che
verrà ripresa anche in seguito

\begin{definition}
    Un polinomio $a_n x^n + \ldots + a_0 \in \ZZx$ si dice
    \textbf{primitivo} se $\MCD(a_n, \ldots, a_0)=1$.
\end{definition}

\begin{theorem}
    \label{th:proiezione_irriducibilità}
    Sia $p$ un primo. Sia $f(x) = a_n x^n + \ldots \in \ZZx$
    primitivo. Se $p \nmid a_n$ e
    $\hatpip(f(x))$ è irriducibile in $\ZZpx$, allora anche $f(x)$ lo
    è in $\ZZx$.
\end{theorem}

\begin{proof}
    Si dimostra la tesi contronominalmente. Sia $f(x) =
        a_nx^n + \ldots \in \ZZ[x]$ primitivo e riducibile, con
    $p \nmid a_n$. Dal momento che $f(x)$ è riducibile, esistono
    $g(x)$, $h(x)$ non invertibili tali che $f(x)=g(x)h(x)$. \\

    Si dimostra che $\deg g(x) \geq 1$. Se infatti fosse nullo,
    $g(x)$ dovrebbe o essere uguale a $\pm 1$ -- assurdo, dal
    momento che $g(x)$ non è invertibile, \Lightning{} -- o
    essere una costante non invertibile. Tuttavia, nell'ultimo
    caso, risulterebbe che $f(x)$ non è primitivo, poiché
    $g(x)$ dividerebbe ogni coefficiente del polinomio.
    Analogamente anche $\deg h(x) \geq 1$. \\

    Si consideri ora $\hatpip(f(x))=\hatpip(g(x))\hatpip(h(x))$.
    Dal momento che $p \nmid a_n$, il grado di $f(x)$ rimane costante
    sotto l'operazione di omomorfismo, ossia $\deg \hatpip(f(x)) =
        \deg f(x)$. \\

    Inoltre, poiché nessuno dei fattori di $f(x)$ è nullo, $\deg f(x) = \deg g(x) +
        \deg h(x)$. Da questa considerazione si deduce che anche i
    gradi di $g(x)$ e $h(x)$ non devono calare, altrimenti si
    avrebbe che $\deg \hatpip(f(x)) < \deg f(x)$, \Lightning{}.
    Allora $\deg \hatpip(g(x)) = \deg g(x) \geq 1$,
    $\deg \hatpip(h(x)) = \deg h(x) \geq 1$. \\

    Poiché $\deg \hatpip(g(x))$ e $\deg \hatpip(h(x))$ sono
    dunque entrambi non nulli, $\hatpip(g(x))$ e $\hatpip(h(x))$
    non sono invertibili\footnote{Si ricorda che $\ZZpx$
        è un anello euclideo. Pertanto, non avere lo stesso grado
        dell'unità equivale a non essere invertibili.}. Quindi
    $f(x)$ è prodotto di non invertibili, ed è dunque riducibile.

\end{proof}

\begin{theorem}[\textit{Criterio di Eisenstein}]
    \label{th:eisenstein}
    Sia $p$ un primo.
    Sia $f(x) = a_n x^n + \ldots + a_0 \in \ZZx$ primitivo tale che:

    \begin{enumerate}[ (1)]
        \item $p \nmid a_n$,
        \item $p \mid a_i$, $\forall i \neq n$,
        \item $p^2 \nmid a_0$.
    \end{enumerate}

    Allora $f(x)$ è irriducibile in $\ZZx$.
\end{theorem}

\begin{proof}
    Si ponga $f(x)$ riducibile e sia pertanto $f(x)=g(x)h(x)$ con
    $g(x)$ e $h(x)$ non invertibili. Analogamente a come visto
    per il \textit{Teorema \ref{th:proiezione_irriducibilità}}, si
    desume che $\deg g(x)$, $\deg h(x) \geq 1$. \\

    Si applica l'omomorfismo di proiezione in $\ZZpx$:

    \[ \hatpip(f(x))=\underbrace{[a_n]_p}_{\neq 0} x_n, \]

    da cui si deduce che $\deg \hatpip(f(x)) = \deg f(x)$. \\

    Dal momento che $\hatpip(f(x))=\hatpip(g(x))\hatpip(h(x))$ e
    che $\ZZpx$, in quanto campo, è un dominio,
    necessariamente sia $\hatpip(g(x))$ che $\hatpip(h(x))$
    sono dei monomi. \\

    Inoltre, sempre in modo analogo a come visto per il \textit{Teorema
        \ref{th:proiezione_irriducibilità}}, sia $\deg \hatpip(g(x))$
    che $\deg \hatpip(h(x))$ sono maggiori o uguali ad $1$. \\

    Combinando questo risultato col fatto che questi due fattori
    sono monomi, si desume che
    $\hatpip(g(x))$ e $\hatpip(h(x))$ sono monomi di grado positivo.
    Quindi $p$ deve dividere entrambi i termini noti di $g(x)$ e
    $h(x)$, e in particolare $p^2$ deve dividere il loro prodotto,
    ossia $a_0$. Tuttavia questo è un assurdo, \Lightning{}.
\end{proof}

\begin{remark*}
    Si consideri $x^k-2$, per $k \geq 1$.
    Per il \nameref{th:eisenstein},
    considerando come primo $p=2$, si verifica che
    $x^k-2$ è sempre irriducibile. Pertanto, per ogni
    grado di un polinomio esiste almeno un irriducibile --
    a differenza di come invece avviene in $\RRx$ o in $\CCx$.
\end{remark*}

\begin{theorem}
    Sia $f(x) \in \ZZx$ primitivo e sia $a \in \ZZ$. Allora $f(x)$ è
    irriducibile se e solo se $f(x+a)$ è irriducibile.
\end{theorem}

\begin{proof}
    Si dimostra una sola implicazione, dal momento che l'implicazione
    contraria consegue dalle stesse considerazioni poste
    studiando prima $f(x+a)$ e poi $f(x)$. \\

    Sia $f(x)=a(x)b(x)$ riducibile, con $a(x)$, $b(x) \in \ZZx$ non
    invertibili. Come già visto per il \textit{Teorema
        \ref{th:proiezione_irriducibilità}}, $\deg a(x)$, $\deg b(x) \geq 1$. \\

    Allora chiaramente $f(x+a)=g(x+a)h(x+a)$, con $\deg g(x+a) =
        \deg g(x) \geq 1$, $\deg h(x+a) = \deg h(x) \geq 1$. Pertanto
    $f(x+a)$ continua a essere riducibile, da cui la tesi.
\end{proof}

\begin{example}
    Si consideri $f(x) = x^{p-1}+\ldots+x^2+x+1 \in \ZZx$, dove
    tutti i coefficienti del polinomio sono $1$. Si verifica che:

    \[ f(x+1)=\frac{(x+1)^p-1}x = p+\binom{p}{2}x+\ldots+x^{p-1}. \]

    Allora, per il \nameref{th:eisenstein} con $p$, $f(x+1)$ è
    irriducibile. Pertanto anche $f(x)$ lo è.
\end{example}

\section{Alcuni irriducibili di \texorpdfstring{$\ZZ_2[x]$}{Z\_2[x]}}

Tra tutti gli anelli $\ZZpx$, $\ZZ_2[x]$ ricopre sicuramente
un ruolo fondamentale, dal momento che è il meno costoso
computazionalmente da analizzare, dacché $\ZZ_2$ consta
di soli due elementi. Pertanto si computano adesso gli
irriducibili di $\ZZ_2[x]$ fino al quarto grado incluso, a meno
di associati. \\

Sicuramente $x$ e $x+1$ sono irriducibili, dal momento che sono di
primo grado. I polinomi di secondo grado devono dunque essere
prodotto di questi polinomi, e pertanto devono avere o $0$ o
$1$ come radice: si verifica quindi che $x^2+x+1$ è l'unico
polinomio di secondo grado irriducibile. \\

Per il terzo grado vale ancora lo stesso principio, per cui
$x^3+x^2+1$ e $x^3+x+1$ sono gli unici irriducibili di tale grado.
Infine, per il quarto grado, i polinomi riducibili soddisfano
una qualsiasi delle seguenti proprietà:

\begin{itemize}
    \item $0$ e $1$ sono radici del polinomio,
    \item il polinomio è prodotto di due polinomi irriducibili di
          secondo grado.
\end{itemize}

Si escludono pertanto dagli irriducibili i polinomi non omogenei --
che hanno sicuramente $0$ come radice --, e i polinomi con $1$ come
radice, ossia $x^4+x^3+x+1$,\ \
$x^4+x^3+x^2+1$, e $x^4+x^2+x+1$. Si esclude anche
$(x^2+x+1)^2 = x^4+x^2+1$. Pertanto gli unici irriducibili di
grado quattro sono $x^4+x^3+x^2+x+1$,\ \ $x^4+x^3+1$,\ \  $x^4+x+1$. \\

Tutti questi irriducibili sono raccolti nella seguente tabella:

\begin{itemize}
    \item (grado 1) $x$, $x+1$,
    \item (grado 2) $x^2+x+1$,
    \item (grado 3) $x^3+x^2+1$, $x^3+x+1$,
    \item (grado 4) $x^4+x^3+x^2+x+1$,\ \ $x^4+x^3+1$,\ \ $x^4+x+1$.
\end{itemize}

\begin{example}
    Il polinomio $51x^3+11x^2+1 \in \ZZx$ è primitivo dal momento
    che $\MCD(51, 11, 1)=1$. Inoltre, poiché $\hatpi_2(51x^3+11x^2+1)=
        x^3+x+1$ è irriducibile, si deduce che anche $51x^3+11x^2+1$ lo
    è per il \textit{Teorema \ref{th:proiezione_irriducibilità}}.
\end{example}

\section{Teorema delle radici razionali e lemma di Gauss}

Si enunciano in questa sezione i teoremi più importanti per
lo studio dell'irriducibilità dei polinomi in $\QQx$ e
in $\ZZx$, a partire dai due teoremi più importanti: il
classico \nameref{th:radici_razionali} e il \nameref{th:lemma_gauss},
che si pone da ponte tra l'analisi dell'irriducibilità in $\ZZx$ e
quella in $\QQx$.

\begin{theorem}[\textit{Teorema delle radici razionali}]
    \label{th:radici_razionali}
    Sia $f(x) = a_n x^n + \ldots + a_0 \in \ZZx$. Abbia $f(x)$
    una radice razionale. Allora, detta tale radice $\frac{p}{q}$,  già ridotta ai minimi termini, questa è tale che:

    \begin{enumerate}[ (i.)]
        \item $p \mid a_0$,
        \item $q \mid a_n$.
    \end{enumerate}
\end{theorem}

\begin{proof}
    Poiché $\frac{p}{q}$ è radice, $f\left(\frac{p}{q}\right)=0$, e
    quindi si ricava che:

    \[ a_n \left( \frac{p}{q} \right)^n + \ldots + a_0 = 0 \implies
        a_n p^n = -q( \ldots + a_0 q^{n-1}). \]

    \vskip 0.1in

    Quindi $q \mid a_n p^n$. Dal momento che $\MCD(p, q)=1$, si
    deduce che $q \mid a_n$. \\

    Analogamente si ricava che:

    \[ a_0 q^n = -p(a_n p^{n-1} + \ldots). \]

    \vskip 0.1in

    Pertanto, per lo stesso motivo espresso in precedenza,
    $p \mid a_0$, da cui la tesi.
\end{proof}

\begin{theorem}[\textit{Lemma di Gauss}]
    \label{th:lemma_gauss}
    Il prodotto di due polinomi primitivi in $\ZZx$ è anch'esso primitivo.
\end{theorem}

\begin{proof}
    Siano $g(x) = a_m x^m + \ldots + a_0$ e $h(x) = b^n x^n + \ldots + b_0$ due polinomi primitivi in $\ZZx$. Si assuma che $f(x)=g(x)h(x)$
    non sia primitivo. Allora esiste un $p$ primo che divide tutti i
    coefficienti di $f(x)$. \\

    Siano $a_s$ e $b_t$ i più piccoli coefficienti non divisibili
    da $p$ dei rispettivi polinomi. Questi sicuramente esistono,
    altrimenti $p$ dividerebbe tutti i coefficienti, e quindi
    o $g(x)$ o $h(x)$ non sarebbe primitivo, \Lightning{}. \\

    Si consideri il coefficiente di $x^{s+t}$ di $f(x)$:

    \[c_{s+t} = \sum_{j+k=s+t} a_j b_k = \underbrace{a_0 b_{s+t} + a_1 b_{s+t-1} + \ldots}_{\equiv \, 0 \pmod p} + a_s b_t + \underbrace{a_{s+1}b_{t-1} + \ldots}_{\equiv \, 0 \pmod p},\]

    dal momento che $p \mid c_{s+t}$, si deduce che $p$ deve dividere
    anche $a_sb_t$, ossia uno tra $a_s$ e $b_t$, che è assurdo, \Lightning{}. Quindi $f(x)$ è primitivo.

\end{proof}

\begin{theorem}[\textit{Secondo lemma di Gauss}]
    \label{th:lemma_gauss_2}
    Sia $f(x) \in \ZZx$. Allora $f(x)$ è irriducibile in $\ZZx$
    se e solo se $f(x)$ è irriducibile in $\QQx$ ed è primitivo.
\end{theorem}

\begin{proof} Si dimostrano le due implicazioni separatamente. \\

    ($\implies$)\; Si dimostra l'implicazione contronominalmente,
    ossia mostrando che se $f(x)$ non è primitivo o se è
    riducibile in $\QQx$, allora $f(x)$ è riducibile in $\ZZx$. \\

    Se $f(x)$ non è primitivo, allora
    $f(x)$ è riducibile in $\ZZx$. Sia quindi $f(x)$ primitivo
    e riducibile in $\QQx$, con $f(x)=g(x)h(x)$,
    $g(x)$, $h(x) \in \QQx \setminus \QQx^*$. \\

    Si descrivano $g(x)$ e $h(x)$ nel seguente modo:

    \[ g(x)=\frac{p_m}{q_m} x^m + \ldots + \frac{p_0}{q_0}, \quad \MCD(p_i, q_i)=1 \; \forall 0 \leq i \leq m, \]

    \[ h(x)=\frac{s_n}{t_n} x^n + \ldots + \frac{s_0}{t_0}, \quad
        \MCD(s_i, t_i)=1 \; \forall 0 \leq i \leq n. \]

    \vskip 0.1in

    Si definiscano inoltre le seguenti costanti:

    \[ \alpha = \frac{\mcm(q_m, \ldots, q_0)}{\MCD(p_m, \ldots, p_0)}, \quad \beta = \frac{\mcm(t_n, \ldots, t_0)}{\MCD(s_n, \ldots, s_0)}. \]

    \vskip 0.1in

    Si verifica che sia $\hat{g}(x)=\alpha g(x)$ che
    $\hat{h}(x)=\beta h(x)$ appartengono a $\ZZx$ e che entrambi
    sono primitivi. Pertanto $\hat{g}(x) \hat{h}(x) \in \ZZx$. \\

    Si descriva $f(x)$ nel seguente modo:

    \[ f(x)=a_k x^k + \ldots + a_0, \quad \MCD(a_k,\ldots,a_0)=1. \]

    \vskip 0.1in

    Sia $\alpha \beta = \frac{p}{q}$ con $\MCD(p,q)=1$, allora:

    \[\hat{g}(x) \hat{h}(x) = \alpha \beta f(x) = \frac{p}{q} (a_k x^k + \ldots + a_0), \]

    da cui, per far sì che $\hat{g}(x) \hat{h}(x)$ appartenga
    a $\ZZx$, $q$ deve necessariamente dividere tutti i
    coefficienti di $f(x)$. Tuttavia $f(x)$ è primitivo, e quindi
    $q=\pm 1$. Pertanto $\alpha \beta = \pm p \in \ZZ$. \\

    Infine, per il \nameref{th:lemma_gauss}, $\alpha \beta f(x)$
    è primitivo, da cui $\alpha \beta = \pm 1$. Quindi
    $f(x) = \pm \hat{g}(x) \hat{h}(x)$ è riducibile. \\

    ($\,\Longleftarrow\,\,$)\; Se $f(x)$ è irriducibile in $\QQx$
    ed è primitivo, sicuramente $f(x)$ è irriducibile anche in
    $\ZZx$. Infatti, se esiste una fattorizzazione in
    irriducibili in $\ZZx$, essa non include alcuna costante
    moltiplicativa dal momento che $f(x)$ è primitivo, e quindi
    esisterebbe una fattorizzazione in irriducibili anche in $\QQx$.
\end{proof}

\end{document}

\newpage
\thispagestyle{empty}
~\newpage

\include{6. Proprietà dei polinomi di K[x] e delle estensioni algebriche}

\newpage
\thispagestyle{empty}
~\newpage

\section{Estensioni algebriche di $\KK$}

\subsection{Morfismi di valutazione, elementi algebrici e trascendenti}

Si definisce adesso il concetto di \textit{omomorfismo di
    valutazione}, che impiegheremo successivamente nello
studio dei quozienti $\KKx/(f(x))$ e dei cosiddetti
\textit{elementi algebrici} (o \textit{trascendenti}).

\begin{definition}
    Sia $B$ un anello commutativo, e sia $A \subseteq B$ un suo
    sottoanello. Si definisce \textbf{omomorfismo di valutazione} di
    $\alpha \in B$ in $A$ l'omomorfismo:

    \[ \valalpha : A[x] \to B, \, f(x) \mapsto f(\alpha). \]
\end{definition}

\begin{remark*}
    L'omomorfismo di valutazione è effettivamente un omomorfismo
    di anelli. Innanzitutto $\valalpha(1)=1$. Inoltre vale
    la linearità:

    \begin{multline*}
        \valalpha(f(x))+\valalpha(g(x))=f(\alpha)+g(\alpha)=
        (f+g)(\alpha)=\valalpha((f+g)(x))=\\=\valalpha(f(x)+g(x)),
    \end{multline*}

    così come la moltiplicatività:

    \begin{multline*}
        \valalpha(f(x))\valalpha(g(x))=f(\alpha)g(\alpha)=
        (fg)(\alpha)=\valalpha((fg)(x))=\valalpha(f(x)g(x)).
    \end{multline*}

    \vskip 0.1in

\end{remark*}

Si evidenziano adesso le principali proprietà di tale
omomorfismo.

\begin{proposition}
    \label{prop:imm_valalpha}
    $\Imm \valalpha = A[\alpha]$
\end{proposition}

\begin{proof} Sicuramente $\Imm \valalpha \subseteq A[\alpha]$,
    dacché ogni immagine di $\valalpha$ è una valutazione di un
    polinomio a coefficienti in $A$ in $\alpha$. \\

    Sia dunque $a=a_n \alpha^n + \ldots + a_0 \in A[\alpha]$. Allora
    $\valalpha(a_n x^n + \ldots + a_0) = a$. Pertanto $a \in \Imm
        \valalpha$, da cui $A[\alpha] \in \Imm \valalpha$. \\

    Poiché vale la doppia inclusione, si desume che $\Imm \valalpha =
        A[\alpha]$.
\end{proof}

Prima di applicare il \textit{Primo teorema d'isomorfismo}, si
distinguono due importanti casi, sui quali si baseranno le
definizioni di \textit{elemento algebrico} e di
\textit{elemento trascendente}.

\begin{definition}
    Sia $\alpha \in B$. Se $\Ker \valalpha = (0)$, allora si
    dice che $\alpha$ è un \textbf{elemento trascendente} di
    $B$ su $A$.
\end{definition}

\begin{remark*}
    Equivalentemente, se $\alpha \in B$ è trascendente su $A$,
    significa che non vi è alcun polinomio non nullo in $A[x]$ che ha $\alpha$
    come soluzione.
\end{remark*}

\begin{example}
    Per esempio, il numero di Nepero-Eulero $e$ è trascendente su $\QQx$\footnote{Per una dimostrazione di questo fatto, si
        guardi a \cite[pp.~234-237]{herstein2010algebra}}. Quindi
    $\Ker \varphi_e = (0)$, e dunque, dal \textit{Primo teorema di
        isomorfismo}, vale che:

    \[ \QQx \cong \QQx/(0) \cong \QQ[e]. \]
\end{example}

Possiamo generalizzare questo esempio nel seguente teorema.

\begin{theorem}
    \label{th:isomorfismo_trascendente}
    Sia $B$ un campo e sia $A \subseteq B$ un suo sottoanello.
    Se $\alpha \in B$ è trascendente su $A$, allora vale
    la seguente relazione:

    \[ A[x] \cong A[\alpha]. \]
\end{theorem}

\begin{proof}
    Si consideri l'omomorfismo $\valalpha$. Dacché $\alpha$ è
    trascendente, $\Ker \valalpha  = (0)$. Allora, combinando
    il \textit{Primo teorema di isomorfismo} con la
    \textit{Proposizione \ref{prop:imm_valalpha}}, si ottiene
    proprio $A[x] \cong A[x]/(0) \cong A[\alpha]$, ossia la tesi.
\end{proof}

\begin{definition}
    Sia $\alpha \in B$. Se $\Ker \valalpha \neq (0)$, allora si
    dice che $\alpha$ è un \textbf{elemento algebrico} di
    $B$ su $A$, mentre il generatore monico\footnote{Vi potrebbero
        essere infatti più generatori di $\Ker \valalpha$, sebbene
        tutti associati tra loro. L'attributo \textit{monico} garantisce
        così l'unicità del polinomio minimo.} non nullo di $\Ker \valalpha$ si
    dice \textbf{polinomio minimo} di $\alpha$ su $A$. Il grado
    di tale polinomio minimo è detto \textbf{grado di} $\alpha$.
\end{definition}

\begin{remark*}
    Equivalentemente, se $\alpha \in B$ è trascendente su $A$,
    significa che esiste un polinomio non nullo in $A[x]$ che ha $\alpha$ come
    soluzione. In particolare, ogni polinomio in $A[x]$ che ha
    $\alpha$ come soluzione è un multiplo del suo polinomio
    minimo su $A$.
\end{remark*}

\begin{example}
    Sia $\alpha \in A$. Allora $\alpha$ è banalmente un elemento
    algebrico su $A$, il cui polinomio minimo è $x-\alpha$. Vale
    dunque che $\Ker \valalpha = (x-\alpha)$, da cui, secondo
    il \textit{Primo teorema di isomorfismo}, si ricava che:

    \[ A[x]/(x-\alpha) \cong A[\alpha] \cong A. \]
\end{example}

\begin{example}
    $i \in \CC$ è un elemento algebrico su $\RR$. Infatti, si
    consideri $\varphi_i$: poiché $i$ è soluzione di $x^2+1$,
    si ha che $x^2+1 \in \Ker \varphi_i$, che è quindi non vuoto. \\

    Inoltre, dal momento che $x^2+1$ è irriducibile in $\RR[x]$,
    esso è generatore di
    $\Ker \varphi_i$. Inoltre, poiché monico, è anche il
    polinomio minimo di $i$ su $\RR$. \\

    Allora, poiché dalla \textit{Proposizione
        \ref{prop:imm_valalpha}} $\Imm \varphi_i = \RR[i]$, si deduce dal \textit{Primo teorema di isomorfismo} che:

    \[ \RRx/(x^2+1) \cong \RR[i] \cong \CC. \]
\end{example}

Ancora una volta possiamo generalizzare questo esempio con il
seguente teorema.

\begin{theorem}
    \label{th:isomorfismo_algebrico}
    Sia $B$ un campo e sia $A \subseteq B$ un suo sottoanello.
    Se $\alpha \in B$ è algebrico su $A$, allora, detto
    $f(x)$ il polinomio minimo di $\alpha$, vale
    la seguente relazione:

    \[ A[x]/(f(x)) \cong A[\alpha]. \]

\end{theorem}

\begin{proof}
    Si consideri l'omomorfismo $\valalpha$. Dacché $\Ker \valalpha
        = (f(x))$ per definizione di polinomio minimo, combinando
    il \textit{Primo teorema di isomorfismo} con la
    \textit{Proposizione \ref{prop:imm_valalpha}}, si ottiene
    proprio $A[x]/(f(x)) \cong A[\alpha]$, ossia la tesi.
\end{proof}

\begin{definition}
    Sia $B$ un campo e sia $A \subseteq B$ un suo sottoanello. Allora,
    dato $\alpha \in B$,
    si definisce con la notazione $A(\alpha)$ il
    sottocampo di $B$ che contiene $A$ e $\alpha$ che
    sia minimale rispetto all'inclusione.
\end{definition}

\begin{remark*}
    Le notazioni $\KK(\alpha, \beta)$ e $\KK(\alpha)(\beta)$ sono equivalenti.
\end{remark*}

\begin{proposition}
    Sia $B$ un campo e sia $A \subseteq B$ un suo sottoanello.
    Se $\alpha \in B$ è algebrico su $A$, allora $A(\alpha)=A[\alpha]$.
\end{proposition}

\begin{proof}
    Se $\alpha$ è algebrico, allora $\Ker \valalpha = (f(x)) \neq (0)$,
    dove $f(x) \in A[x]$ è irriducibile. Pertanto, per
    il \textit{Teorema \ref{th:campo_quoziente_irriducibile}},
    $A[x]/(f(x))$ è un campo. \\

    Dunque dal \textit{Teorema \ref{th:isomorfismo_algebrico}} si
    ricava che:

    \[ A[x]/(f(x)) \cong A[\alpha]. \]

    \vskip 0.1in

    Pertanto $A[\alpha]$ è un campo. Dacché $A[\alpha] \subseteq A(\alpha)$ e $A(\alpha)$ è minimale rispetto all'inclusione,
    si deduce che $A[\alpha]=A(\alpha)$, ossia la tesi.
\end{proof}

\begin{remark*}
    Il teorema che è stato appena enunciato non vale per
    gli elementi trascendenti. Infatti, $A[\alpha]$ sarebbe
    isomorfo a $A[x]$, che non è un campo. Al contrario
    $A(\alpha)$ è un campo, per definizione.
\end{remark*}

\begin{proposition}
    Sia $B$ un campo e sia $A \subseteq B$ un suo sottoanello.
    Se $\alpha$, $\beta \in B$ sono algebrici su $A$ e condividono
    lo stesso polinomio minimo, allora $A[\alpha] \cong A[\beta]$.
\end{proposition}

\begin{proof}
    Sia $f(x)$ il polinomio minimo di $\alpha$ e $\beta$.
    Dal \textit{Primo teorema di isomorfismo} e dalla
    \textit{Proposizione \ref{prop:imm_valalpha}} si
    desume che $A[x]/(f(x)) \cong A[\alpha]$. Analogamente
    si ricava che $A[x]/(f(x)) \cong A[\beta]$. Pertanto
    $A[\alpha] \cong A[\beta]$.
\end{proof}

\subsection{Teorema delle torri ed estensioni algebriche}

\begin{definition}
    Siano $A \subseteq B$ campi. Allora si denota come
    $[B : A]$ la dimensione dello spazio vettoriale $B$
    costruito su $A$, ossia $\dim B_A$. Tale dimensione è detta \textbf{grado
        dell'estensione}.
\end{definition}

\begin{theorem}[\textit{Teorema delle torri algebriche}]
    \label{th:torri}
    Siano $A \subseteq B \subseteq C$ campi. Allora:

    \[ [C : A] = [C : B] [B : A]. \]
    \vskip 0.1in
\end{theorem}

\begin{proof}
    Siano $[C : B] = m$ e $[B : A] = n$. Sia
    $\BB_C = (a_1, \ldots, a_m)$ una base
    di $C$ su $B$, e sia $\BB_B = (b_1, \ldots, b_n)$ una
    base di $B$ su $A$. \\

    Si dimostra che la seguente è una base di $C$ su $A$:

    \[\BB_A \BB_B = \{ a_1b_1, \ldots, a_1b_n, \ldots, a_mb_n\}. \]

    \vskip 0.1in

    \ (i) $\BB_C \BB_B$ genera $A$ su $C$. \\

    Sia $c \in C$. Allora si può descrivere $a$ nel seguente
    modo:

    \[c = \sum_{i=1}^m \beta_i a_i, \quad \text{con } \beta_i \in B, \; \forall 1 \leq i \leq m.\]

    A sua volta, allora, si può descrivere ogni $\beta_i$ nel
    seguente modo:

    \[\beta_i = \sum_{j=1}^n \gamma_j^{(i)} b_j, \quad \text{con }
        \gamma_j^{(i)} \in A, \; \forall 1 \leq j \leq n.\]

    \vskip 0.1in

    Combinando le due equazioni, si verifica che $\BB_C \BB_B$ genera $C$ su $A$:

    \[ c = \sum_{i=1}^m \sum_{j=1}^n \gamma_j^{(i)} b_j a_i, \quad \text{con } \gamma_j^{(i)} \in A, \; \forall 1 \leq i \leq m, \, 1 \leq j \leq n. \]

    \vskip 0.1in

    \ (ii) $\BB_C \BB_B$ è linearmente indipendente. \\

    Si consideri l'equazione:

    \[ \sum_{i=1}^m \sum_{j=1}^n \gamma_j^{(i)} b_j a_i = 0, \quad \text{con } \gamma_j^{(i)} \in A, \; \forall 1 \leq i \leq m, \, 1 \leq j \leq n .\]

    Poiché $\BB_C$ è linearmente indipendente, si deduce
    che:

    \[ \sum_{j=1}^n \gamma_j^{(i)} b_j = 0, \; \forall 1 \leq i \leq m. \]

    Tuttavia, $\BB_B$ è a sua volta linearmente indipendente,
    e quindi $\gamma_j^{(i)} = 0$, $\forall i, j$. Dunque
    $\BB_C \BB_B$ è linearmente indipendente. \\

    Dal momento che $\BB_C \BB_B$ è linearmente indipendente e
    genera $C$ su $A$, consegue che essa sia una base di $C$ su
    $A$. Quindi $[C : A] = mn = [C : B][B : A]$, da cui la tesi.
\end{proof}

\begin{definition}
    Siano $A \subseteq B$ campi. Se $[B : A] \neq \infty$, allora
    si dice che $BA$ è un'\textbf{estensione finita} di $A$.
    Altrimenti si dice che $B$ è un'\textbf{estensione infinita}
    di $A$.
\end{definition}

\begin{proposition}
    \label{prop:estensione_finita}
    Siano $A \subseteq B \subseteq C$ campi. Allora, se $C$ è
    un'estensione finita di $A$, anche $B$ lo è. Inoltre
    $C$ è un'estensione finita di $B$.
\end{proposition}

\begin{proof}
    Dal momento che $B$ è un sottospazio dello spazio vettoriale
    $C$ costruito su $A$, e questo ha dimensione finita,
    anche $B$ su $A$ ha dimensione finita. Quindi $[B : A] \neq
        \infty$, e $B$ è dunque un'estensione finita di $A$. \\

    Infine, dacché una base di $C$ su $A$ è un generatore finito
    di $C$ su $B$, si deduce che $[C : B] \neq \infty$, e quindi
    che $C$ è un'estensione finita di $B$.
\end{proof}

\begin{theorem}
    \label{th:estensione_algebrica}
    Siano $A \subseteq B$ campi. Allora $a \in B$ è
    algebrico su $A$ se e solo se $[A(a) : A] \neq \infty$,
    ossia solo se $A(a)$ è un'estensione finita di $A$.
\end{theorem}

\begin{proof} Si dimostrano le due implicazioni separatamente. \\

    ($\implies$)\; Se $a \in B$ è algebrico su $A$, allora
    dal \textit{Teorema \ref{th:isomorfismo_algebrico}} si ricava che:

    \[ A[x]/(f(x)) \cong A[a] \cong A(a). \]

    \vskip 0.1in

    Dacché $A[x]/(f(x))$ ha dimensione finita, anche $A(a)$
    ha dimensione finita, e quindi è un'estensione finita
    di $A$. \\

    ($\,\Longleftarrow\,\,$)\; Sia $A(a)$ un'estensione
    finita di $A$ e sia $[A(a) : A]=m$. Allora $I=(1, a, a^2, \ldots, a^m)$ è linearmente dipendente, dal momento che contiene
    $m+1$ elementi. Quindi esiste una sequenza finita non nulla
    $(\alpha_i)_{i=\,0\to m}$ con elementi in $A$ tale che:

    \[ \alpha_m a^m + \ldots + \alpha_2 a^2 + \alpha_1 a + \alpha_0 = 0. \]

    Quindi $a$ è soluzione del polinomio:

    \[ f(x) = \alpha_m x^m + \ldots + \alpha_2 x^2 + \alpha_1 x + \alpha_0 \in A[x], \]

    \vskip 0.1in

    pertanto $a$ è algebrico su $A$, da cui la tesi.
\end{proof}

\begin{definition}
    Siano $A \subseteq B$ campi. Allora si dice che $B$ è
    un'\textbf{estensione algebrica} di $A$ se ogni elemento
    di $B$ è algebrico su $A$.
\end{definition}

\begin{proposition}
    \label{prop:estensione_finita_algebrica}
    Siano $A \subseteq B$ campi. Se $B$ è un'estensione finita
    di $A$, allora $B$ è una sua estensione algebrica.
\end{proposition}

\begin{proof}
    Sia $\alpha \in B$ e si consideri la catena di campi $A \subseteq A(\alpha)
        \subseteq B$. Dacché $[B : A] \neq \infty$, per la \propref{prop:estensione_finita}
    anche $[A(\alpha) : A] \neq \infty$. Pertanto, dal \thref{th:estensione_algebrica}, $\alpha$ è algebrico. Così tutti gli elementi
    di $B$ sono algebrici in $A$, e dunque, per definizione, $B$ è un'estensione
    algebrica di $A$.
\end{proof}

\begin{theorem}
    \label{th:somma_prodotto_algebrici}
    Siano $A \subseteq B$ campi e siano $\beta_1$, $\beta_2$, $\ldots$, $\beta_n$
    elementi algebrici di $B$ su $A$, con $n \geq 1$.
    Allora $[A(\beta_1, \beta_2, \ldots, \beta_n) : A] \neq \infty$.
\end{theorem}

\begin{proof} Si procede applicando il principio di induzione su $n$. \\

    \ (\textit{passo base}) La tesi è verificata per il \thref{th:estensione_algebrica}. \, \\

    \ (\textit{passo induttivo}) Per l'ipotesi induttiva, si sa che
    $[A(\beta_1, \beta_2, \ldots, \beta_{n-1}) : A] \neq \infty$. \\

    Poiché $\beta_n$ è algebrico su $A$, sin da subito si osserva
    che $[A(\beta_n) : A] \neq \infty$ per il \thref{th:estensione_algebrica}.
    Sia allora $f(x)$ il polinomio minimo di $\beta_n$ appartenente a
    $A[x]$. Esso è un polinomio che ammette $\beta_n$ come radice
    anche in $A(\beta_1, \beta_2, \ldots, \beta_{n-1})[x]$, e quindi
    $\Ker \varphi_{\beta_n} \neq (0)$ ammette un generatore
    $p(x)$, che divide $f(x)$. Si ottiene pertanto la seguente
    disuguaglianza:

    \[ [A(\beta_1, \beta_2, \ldots, \beta_{n-1})(\beta_n) : A(\beta_1, \beta_2, \ldots, \beta_{n-1})] = \deg p(x) \leq
        \deg f(x) = [A(\beta_n) : A].  \]
        
    \vskip 0.1in

    Poiché $[A(\beta_n) : A]$ è finito, anche $[A(\beta_1, \beta_2, \ldots, \beta_{n-1})(\beta_n) : A(\beta_1, \beta_2, \ldots, \beta_{n-1})]$ lo è. \\

    Combinando i due risultati, si ottiene con il \nameref{th:torri} che:

    \begin{multline*}
        [A(\beta_1, \beta_2, \ldots, \beta_n) : A] = [A(\beta_1, \beta_2, \ldots, \beta_{n-1})(\beta_n) : A(\beta_1, \beta_2, \ldots, \beta_{n-1})] \\ \cdot[A(\beta_1, \beta_2, \ldots, \beta_{n-1}) : A] \neq \infty,
    \end{multline*}

    da cui la tesi.
    \, \\
\end{proof}

\begin{corollary}
    \label{cor:estensione_algebrica_due_elementi}
    Siano $A \subseteq B$ campi e siano $\alpha$, $\beta \in B$ elementi
    algebrici su $A$. Allora $A(\alpha, \beta)$ è un'estensione algebrica.
\end{corollary}

\begin{proof}
    Dal \thref{th:somma_prodotto_algebrici} si ricava che $[A(\alpha, \beta) : A] \neq
        \infty$. Quindi $A(\alpha, \beta)$ è un'estensione finita di $A$, ed in quanto
    tale, per la \propref{prop:estensione_finita_algebrica}, essa è algebrica.
\end{proof}

\begin{remark*}
    Esistono estensioni algebriche che hanno grado infinito. Un
    esempio notevole è $\mathcal{A}$, l'insieme dei numeri algebrici di $\CC$
    su $\QQ$. Infatti, si ponga $[\mathcal{A} : \QQ] = n-1 \in \NN$ e si
    consideri $x^n-2$. Dal momento che per il \textit{Criterio di Eisenstein}
    tale polinomio è irriducibile, si ricava che $[\QQ(\nsqrt{n}{2}) : \QQ] = n$. \\

    Poiché $\nsqrt{n}{2}$ è algebrico, si deduce che $\QQ(\nsqrt{n}{2}) \subseteq
        \mathcal{A}$, dal momento che per il \corref{cor:estensione_algebrica_due_elementi} ogni elemento di $\QQ(\nsqrt{n}{2})$ è algebrico su $\QQ$.
    Tuttavia questo è un assurdo dal momento che
    $\QQ(\nsqrt{n}{2})$ ha
    dimensione maggiore di $\mathcal{A}$, di cui è sottospazio vettoriale.
\end{remark*}

\begin{proposition}
    \label{prop:alpha_quadro}
    Siano $A \subseteq B$ campi e sia $\alpha \in B$. Se $[A(\alpha) : A]$
    è dispari, allora $A(\alpha^2)=A(\alpha)$.
\end{proposition}

\begin{proof}
    Innanzitutto, si osserva che $A(\alpha^2) \subseteq A(\alpha)$, ossia
    che $A(\alpha)$ è un'estensione di $A(\alpha^2)$. Grazie a questa
    osservazione è possibile considerare il grado di $A(\alpha)$ su
    $A(\alpha^2)$, ossia $[A(\alpha) : A(\alpha^2)]$. Poiché $\alpha$ è
    radice del polinomio $x^2 - \alpha^2$ in $A(\alpha^2)$, si deduce
    che tale grado è al più $2$. \\
    
    Si applichi il \nameref{th:torri} alla catena di estensioni
    $A \subseteq A(\alpha^2) \subseteq A(\alpha)$:
    
    \[ [A(\alpha) : A] = \underbrace{[A(\alpha) : A(\alpha^2)]}_{\leq 2} [A(\alpha^2) : A]. \]
    
    \vskip 0.1in
    
    Se $[A(\alpha) : A(\alpha^2)]$ fosse $2$, $[A(\alpha) : A]$ sarebbe
    pari, \Lightning{}. Pertanto $[A(\alpha) : A(\alpha^2)] = 1$, da
    cui si ricava che $[A(\alpha) : A] = [A(\alpha^2) : A]$, ossia
    che $A(\alpha^2)$ ha la stessa dimensione di $A(\alpha)$ su $A$. \\
    
    Dal momento che $A(\alpha^2)$ è un sottospazio vettoriale di $A(\alpha)$,
    avere la sua stessa dimensione equivale a coincidere con lo spazio
    stesso. Si conclude allora che $A(\alpha^2) = A(\alpha)$.
\end{proof}

\begin{remark*}
    Si osserva che la \propref{prop:alpha_quadro} si può generalizzare
    facilmente ad un esponente $n$ qualsiasi, finché sia data come ipotesi
    la non divisibilità di $[A(\alpha) : A]$ per nessun numero primo
    minore o uguale di $n$. \\
    
    Si può infatti considerare, per
    la dimostrazione generale, il polinomio $x^n - \alpha^n$, la cui
    esistenza implica che $[A(\alpha) : A(\alpha^n)]$ sia minore
    o uguale di $n$.
\end{remark*}


\begin{theorem}
    Siano $A \subseteq B \subseteq C$ campi. Se $B$ è un'estensione algebrica di $A$
    e $C$ è un'estensione algebrica di $B$, allora $C$ è un'estensione algebrica di
    $A$.
\end{theorem}

\begin{proof}
    Per mostrare che $C$ è un'estensione algebrica di $A$, verificheremo che
    ogni suo elemento è algebrico in $A$. Sia dunque $c \in C$. \\

    Poiché per ipotesi $c$ è algebrico su $B$, esiste un polinomio $f(x) \in B[x]$
    tale che $c$ ne sia radice. Sia $f(x)$ il polinomio minimo di $c$ su $B$,
    descritto come:

    \[ f(x) = b_0 + b_1 x + \ldots + b_n x^n,\quad n = [B(c) : B].\]

    \vskip 0.1in

    Dacché $B$ è un'estensione algebrica di $A$, ogni coefficiente $b_i$ di $f(x)$ è
    algebrico su $A$, ossia $[A(b_i) : A] \neq \infty$. Allora, per il
    \thref{th:somma_prodotto_algebrici}, $[A(b_0, \ldots, b_n) : A] \neq \infty$.
    \\

    Anche $[A(c, b_0, \ldots, b_n) : A(b_0, \ldots, b_n)] \neq \infty$, dal
    momento che $c$ è soluzione di $f(x) \in A(b_0, \ldots, b_n)[x]$. \\

    Allora, per il \nameref{th:torri}, $[A(c, b_0, \ldots, b_n) : A] = [A(c, b_0,
        \ldots, b_n) : A(b_0, \ldots, b_n)][A(b_0, \ldots, b_n) : A] \neq \infty$.
    Quindi $A(c, b_0, \ldots, b_n)$ è un'estensione finita di $A$. \\

    Poiché $A \subseteq A(c) \subseteq A(c, b_0, \ldots, b_n)$ è una
    catena di estensione di campi, per la \propref{prop:estensione_finita},
    $A(c)$ è un'estensione finita di $A$, ed in quanto tale, per
    la \propref{prop:estensione_finita_algebrica}, è anche algebrica. Quindi
    $c$ è algebrico su $A$, da cui la tesi.
\end{proof}

\newpage
\thispagestyle{empty}
~\newpage

\section{Teorema fondamentale dell'Algebra e radici reali in \texorpdfstring{$\QQx$}{Q[x]}}

Si enuncia adesso il \nameref{th:algebra}, senza tuttavia
fornirne una dimostrazione\footnote{Per la dimostrazione si rimanda
    a \cite[pp.~142-143]{di2013algebra}, avvisando della sua
    estrema tecnicità. Una dimostrazione a tema strettamente
    algebrico è dovuta invece al matematico francese Laplace (1749 -- 1827), per la quale
    si rimanda a \cite[pp.~120-122]{Remmert1991}.}.

\begin{theorem}[\textit{Teorema fondamentale dell'Algebra}]
    \label{th:algebra}
    Un polinomio non costante $f(x) \in \CCx$ ammette sempre almeno una radice in
    $\CC$.
\end{theorem}

\begin{corollary}
    Sia $f(x) \in \CCx$ di grado $n\geq1$. Allora $f(x)$ ammette
    esattamente $n$ radici, contate con la giusta molteplicità.
\end{corollary}

\begin{proof}
    Sia $\zeta_1$ una radice complessa di $f(x)$, la cui esistenza
    è garantita dal \nameref{th:algebra}. Si divida $f(x)$ per
    $(x-\zeta_1)$ e se ne prende il quoziente $q_1(x)$, mentre si
    ignori il resto, che
    per la \textit{Proposizione \ref{prop:radice_x_meno_alpha}},
    è nullo. \\

    Si reiteri il procedimento utilizzando $q_1(x)$ al
    posto di $f(x)$ fino a quando il grado del quoziente non è nullo,
    e si chiami infine questo quoziente di grado nullo $\alpha$.
    Infatti, poiché i gradi dei quozienti diminuiscono di $1$ ad
    ogni iterazione, è garantito che l'algoritmo termini esattamente
    dopo $n$ iterazioni. Pertanto, $f(x)$ a priori ha almeno $n$ radici. \\

    In questo modo, numerando le radici, si può scrivere $f(x)$ come:

    \begin{equation}
        \label{eq:fattorizzazione_fx__reali}
        f(x)=\alpha(x-\zeta_1)(x-\zeta_2)\cdots(x-\zeta_n).
    \end{equation}

    \vskip 0.1in

    Dal momento che $x-\zeta_i$ è irriducibile $\forall 1 \leq i \leq n$
    e dacché $\KKx$, in quanto anello euclideo, è un UFD, si dimostra
    che \eqref{eq:fattorizzazione_fx__reali} è l'unica fattorizzazione di
    $f(x)$, a meno di associati. Pertanto $f(x)$ ammette esattamente
    $n$ radici.
\end{proof}

\newpage
\thispagestyle{empty}
~\newpage

\chapter{Introduzione alla teoria dei campi}

\section{La caratteristica di un campo}

Si consideri il seguente omomorfismo:

\[ \psi : \ZZ \to \KK, \]

\vskip 0.1in

completamente determinato dalla condizione $\psi(1) = 1$, dacché
$\ZZ$ è generato da $1$. Si studia innanzitutto il caso in cui
$\Ker \psi = (0)$. In questo caso, $\psi$ è un monomorfismo, e per
il \corref{cor:primo_isomorfismo_iniettivo}, $\ZZ \cong \Imm \psi$. \\

Pertanto, $\KK$ ammetterebbe come sottoanello una copia isomorfa di $\ZZ$.
Inoltre, poiché $\KK$ è un campo, deve anche ammetterne gli inversi, e quindi
ammetterebbe come sottocampo una copia isomorfa di $\QQ$. La seguente
definizione classificherà questi tipi di campo. \\

\begin{definition}
    Si dice che un campo $\KK$ è di \textbf{caratteristica zero} ($\Char \KK = 0$),
    quando $\Ker \psi = (0)$.
\end{definition}

Altrimenti, se $\Ker \psi \neq (0)$, dacché $\ZZ$ è un anello euclideo,
$\Ker \psi$ deve essere monogenerato da un intero $n$, ossia $\Ker \psi = (n)$. \\

Tuttavia non tutti gli interi sono ammissibili. Sia infatti $n$ non primo, allora
$n = ab$ con $a$, $b \neq \pm 1$. Si nota innanzitutto che $\psi(a) \neq 0$,
se infatti fosse nullo, $n$ dovrebbe dividere $a$, impossibile dal momento
che $\card{a} < \card{n}$, \Lightning{}. Analogamente anche $\psi(b) \neq 0$. \\

Se $n$ fosse generatore di $\Ker \psi$ si ricaverebbe allora che:

\[ \underbrace{\psi(a)}_{\neq\,0} \underbrace{\psi(b)}_{\neq\,0} = \psi(n) = 0, \]

\vskip 0.1in

che è assurdo, dal momento che $\KK$, in quanto campo, è anche un dominio.
Quindi $n$ deve essere un numero primo. In particolare, allora, per
il \nameref{th:primo_isomorfismo}, $\ZZp = \ZZ/(p) \cong \Imm \psi$,
ossia $\KK$ contiene una copia isomorfa di $\ZZp$, a cui ci riferiremo
semplicemente con $\FFpp$. \\

Allora, poiché sia $\KK$ che $\FFpp$ sono campi, $\KK$ è uno spazio
vettoriale su $\FFpp$. Si può dunque classificare quest'ultimo tipo di
campi con la seguente definizione:

\begin{definition}
    Si dice che un campo $\KK$ è di \textbf{caratteristica $p$} ($\Char \KK = p$)
    quando $\Ker \psi = (p)$, con $p$ primo.
\end{definition}

\begin{remark*}
    La caratteristica di un campo \textbf{non} distingue i campi finiti
    dai campi infiniti. Esistono infatti campi infiniti di caratteristica
    $p$, come il campo delle funzioni razionali su $\ZZp$:

    \[ \ZZ_p(x) =  \left\{ \frac{f(x)}{g(x)} \mid f(x),\, g(x) \in \ZZpx,\, g(x) \neq 0 \right\}. \]

    \vskip 0.1in

    Infatti $\psi(p) = p \, \psi(1) = 0$.
\end{remark*}

\section{Prime proprietà dei campi di caratteristica \texorpdfstring{$p$}{p}}

Come si è appena visto, un campo $\KK$ di caratteristica $p$ contiene
al suo interno un sottocampo $\FFpp$ isomorfo a $\ZZp$, ed è per questo
uno spazio vettoriale su di esso. A partire da questa informazione si
può dimostrare la seguente proposizione.

\begin{proposition}
    \label{prop:campo_char_p_prodotto_per_p}
    Sia $\KK$ un campo di caratteristica $p$. Allora, per ogni
    elemento $v$ di $\KK$, $pv=0$.
\end{proposition}

\begin{proof}
    Considerando ogni elemento di $\KK$ come vettore e $p$ come
    scalare, si ricava che:

    \[ pv=(\underbrace{1+\ldots+1}_{p\text{ volte}})v=
        (\underbrace{\psi(1)+\ldots+\psi(1)}_{p\text{ volte}})v=
        \psi(p)v=0v=0. \]

\end{proof}

Mentre, partendo da questa proposizione, si può dimostrare il
seguente teorema.

\begin{theorem}[\textit{Teorema del binomio ingenuo}]
    \label{th:binomio_ingenuo}
    Siano $a$ e $b$ elementi di un campo di caratteristica $p$. Allora
    $(a+b)^p = a^p + b^p$.
\end{theorem}

\begin{proof}
    Per dimostrare la tesi si applica la formula del binomio di Newton
    nel seguente modo:

    \[ (a+b)^p = \sum_{i=0}^p \binom{p}{i} a^{p-i}b^p. \]

    \vskip 0.1in

    Tuttavia, dal momento che $p$ è un fattore di tutti i binomiali per
    $1 \leq i \leq p-1$, tutti i termini computati con queste $i$
    sono nulli per la \propref{prop:campo_char_p_prodotto_per_p}.
    Si desume così l'identità della tesi.
\end{proof}

\section{L'omomorfismo di Frobenius}

\begin{definition}
    Dato un campo $\KK$ di caratteristica $p$, si definisce
    \textbf{omomorfismo di Frobenius} per il campo $\KK$
    la funzione:

    \[ \Frob : \KK \to \KK,\, a \mapsto a^p. \]
\end{definition}

\begin{remark*}
    In effetti, l'omomorfismo di Frobenius è un omomorfismo. \\

    Infatti, $\Frob(1) = 1^p = 1$. Inoltre tale funzione
    rispetta la linearità per il \nameref{th:binomio_ingenuo}:

    \[ \Frob(a + b) = (a+b)^p = a^p + b^p = \Frob(a) + \Frob(b), \]

    \vskip 0.1in

    e chiaramente anche la moltiplicatività:

    \[ \Frob(ab) = (ab)^p = a^p b^p = \Frob(a) \Frob(b). \]
\end{remark*}

\begin{proposition}
    \label{prop:frobenius_monomorfismo}
    L'omomorfismo di Frobenius di un campo $\KK$ di caratteristica
    $p$ è un monomorfismo.
\end{proposition}

\begin{proof}
    Si prenda in considerazione $\Ker \Frob$. Esso è sicuramente
    un ideale diverso da $\KK$, dacché $1 \notin \Ker \Frob$.
    Tuttavia, se $\Ker \Frob \neq (0)$, $\Ker \Frob$, dal
    momento che $\KK$, in quanto campo, è un anello euclideo,
    e quindi un PID, è monogenerato da un invertibile. \\

    Se però così fosse, $\Ker \Frob$ coinciderebbe con il
    campo $\KK$ stesso, \Lightning{}. Quindi $\Ker \Frob = (0)$,
    da cui la tesi.
\end{proof}

\begin{proposition}
    Sia $\KK$ un campo finito di caratteristica $p$. Allora
    l'omomorfismo di Frobenius è un automorfismo.
\end{proposition}

\begin{proof}
    Dalla \propref{prop:frobenius_monomorfismo} è noto che
    $\Frob$ sia già un monomorfismo. Dal momento che
    il dominio e il codominio sono lo stesso e constano
    entrambi dunque di un numero finito di elementi,
    se $\Frob$ non fosse surgettivo, vi sarebbe un elemento
    di $\KK$ a cui non è associato nessun elemento di $\KK$
    mediante $\Frob$. \\

    Per il principio dei cassetti, allora, spartendo
    $\card{\KK}$ elementi in $\card{\KK}-1$ elementi,
    vi sarebbe almeno un elemento dell'immagine a cui
    sarebbero associati due elementi del dominio. Tuttavia
    questo è assurdo dal momento che $\Frob$ è un
    monomorfismo. Quindi $\Frob$ è un epimorfismo. \\

    Dacché $\Frob$ è contemporaneamente un endomorfismo,
    un monomorfismo e un epimorfismo, è allora anche
    un automorfismo.
\end{proof}

\begin{proposition}
    \label{prop:punti_fissi_frobenius_campo}
    Sia $\KK$ un campo di caratteristica $p$ e si
    definisca l'insieme dei punti fissi del suo
    omomorfismo di Frobenius:

    \[ \Fix(\Frobexp^n) = \{ a \in \KK \mid \Frobexp^n(a) = a \} .\]

    \vskip 0.1in

    Allora $\Fix(\Frobexp^n)$ è un sottocampo di $\KK$.
\end{proposition}

\begin{proof}
    Affinché $\Fix(\Frobexp^n)$ sia un sottocampo di $\KK$,
    la sua somma e la sua moltiplicazione devono essere ben
    definite, e ogni suo elemento deve ammettere un inverso
    sia additivo che moltiplicativo. \\

    Siano allora $a$, $b \in \Fix(\Frobexp^n)$.
    $\Frobexp^n$ è un omomorfismo, in quanto è composizione
    di omomorfismi (in particolare, dello stesso omomorfismo
    $\Frobexp$). Sfruttando le proprietà
    degli omomorfismi si dimostra dunque
    che $a+b \in \Fix(\Frobexp^n)$:

    \[ \Frobexp^n(a+b) = \Frobexp^n(a) + \Frobexp^n(b) = a+b, \]

    \vskip 0.1in

    e che $ab \in \Fix(\Frobexp^n)$:

    \[ \Frobexp^n(ab) = \Frobexp^n(a)\Frobexp^n(b) = ab. \]

    \vskip 0.1in

    Analogamente si dimostra che $-a \in \Fix(\Frobexp^n)$:

    \[ \Frobexp^n(-a) = -\Frobexp^n(a) = -a, \]

    \vskip 0.1in

    e che $a\inv \in \Fix(\Frobexp^n)$:

    \[ \Frobexp^n(a\inv) = \Frobexp^n(a)\inv = a\inv. \]

\end{proof}


\section{Classificazione dei campi finiti}

\begin{theorem}
    Ogni campo finito $\KK$ di caratteristica $p$ consta
    di $p^n$ elementi, con $n \in \NN^+$.
\end{theorem}

\begin{proof}
    Come già detto precedentemente, $\KK$ è uno
    spazio vettoriale su una copia isomorfa di $\ZZp$,
    $\FFpp$. \\

    Si consideri allora il grado $[\KK : \FFpp]$. Sicuramente
    questo grado non è infinito, dal momento che $\KK$ non
    ha infiniti elementi. Quindi $[\KK : \FFpp] = n \in \NN$. \\

    Sia dunque $(k_1, k_2, \ldots, k_n)$ una base di $\KK$
    su $\FFpp$. Ogni elemento $a$ di $\KK$ si potrà dunque scrivere
    come:

    \[ a = \alpha_1 k_1 + \ldots + \alpha_n k_n, \quad \alpha_1, \ldots, \alpha_n \in \FFpp,\]

    \vskip 0.1in

    e dunque vi saranno in totale $p^n$
    elementi, dove ogni $p$ è contato dal numero di elementi che è
    possibile associare ad ogni coefficiente, ossia $\card{\FFpp} = p$,
    per il numero di elementi appartenenti alla base, ossia $[\KK : \FFpp] =
        n$, da cui la tesi.
\end{proof}

\begin{theorem}
    Per ogni $n \in \NN^+$ e per ogni numero primo $p$ esiste un
    campo finito con $p^n$ elementi.
\end{theorem}

\begin{proof}
    Si consideri il polinomio $x^{p^n}-x$ su $\ZZp$ e un suo
    campo di spezzamento $A$. $\Fix(\Frobexp^n)$, per
    la \propref{prop:punti_fissi_frobenius_campo}, è
    un sottocampo, e
    contiene esattamente le radici di $x^{p^n}-x$, che
    in $A$ si spezza in fattori lineari, per definizione. \\

    La derivata di $x^{p^n}-x$ è $p^n x^{p^n - 1}-1 \equiv -1$,
    dacché $A$ è uno spazio vettoriale su $\ZZp$, e pertanto
    vale ancora la \propref{prop:campo_char_p_prodotto_per_p}.
    Dal momento che $-1$ e $x^{p^n}-x$ non hanno fattori lineari
    in comune, per il \textit{Criterio della derivata},
    $x^{p^n}-x$ non ammette radici multiple. \\

    Allora $\Fix(\Frobexp^n)$ è un campo con $p^n$ elementi,
    ossia tutte le radici di $x^{p^n}-x$ (e coincide quindi
    con il campo di spezzamento $A$), da cui la tesi.
\end{proof}


\newpage
\thispagestyle{empty}
~\newpage

\section{Teoremi rilevanti sui campi finiti}

\subsection{Campo di spezzamento di un irriducibile in \texorpdfstring{$\FFpp$}{F\_p}}

\begin{theorem}
    Sia $f(x)$ un polinomio irriducibile in $\FFpp$ e sia
    $n$ il suo grado. Allora $\FFpn$ è il suo campo
    di spezzamento.
\end{theorem}

\begin{proof}
    Dacché $f(x)$ è irriducibile, $\FFpp/((f(x))$ è un
    campo con $p^n$ elementi, ed è quindi isomorfo
    a $\FFpn$. \\

    Sia $\alpha = x + (f(x))$ una radice di $f(x)$
    in $\FFpn$. Dal momento che $f(x)$ è irriducibile in
    $\FFpp$, esso è il polinomio minimo di $\alpha$. Tuttavia,
    poiché $\alpha \in \FFpn$, $\alpha$ è anche radice
    di $x^{p^n}-x$. Pertanto si deduce che $f(x)$ divide
    $x^{p^n}-x$. \\

    Dunque, poiché $x^{p^n}-x$ in $\FFpn$ è prodotto di
    fattori lineari, tutte le radici di $f(x)$ sono già
    in $\FFpn$. \\

    Inoltre, $\FFpn$ è il più piccolo sottocampo contenente
    $\alpha$, dacché $\FFpn \cong \FFpp/(f(x)) \cong \FFpp(\alpha)$.
    Quindi si deduce che $\FFpn$ è un campo di spezzamento per
    $f(x)$, ossia la tesi.
\end{proof}

\begin{lemma}
    \label{lem:frobexp}
    Sia $f(x)$ un irriducibile di grado $n$ su $\FFpp[x]$ e sia $\alpha$
    una sua radice in $\FFpn$. Allora $f(\Frobexp^k(\alpha))=0$, $\forall k \geq 0$
    \footnote{$\Frob$ è l'omomorfismo di Frobenius, definito come $\Frob : \FFpp \to \FFpp$,
        $a \mapsto a^p$.}.
\end{lemma}

\begin{proof} Sia $f(x) = a_n x^n + \ldots + a_0$ a coefficienti in $\FFpp$.
    Si dimostra la tesi applicando il principio di induzione su $k$. \\

    \ (\textit{passo base})\; $f(\Frobexp^0(\alpha))=f(\alpha)=0$. \\

    \ (\textit{passo induttivo})\; Per l'ipotesi induttiva, $f(\Frobexp^{k-1}(\alpha))=0$.
    Allora, si verifica algebricamente che:

    \begin{multline*}
        f(\Frobexp^k(\alpha)) = a_n (\Frobexp^k(\alpha))^n + \ldots + a_0 =
        \Frob(a_n) \Frob((\Frobexp^{k-1}(\alpha))^n) + \ldots + \Frob(a_0) = \\
        \Frob(f(\Frobexp^{k-1}(\alpha))) = \Frob(0) = 0,
    \end{multline*}

    \vskip 0.1in

    dove si è usato che $\Frob(a_i) = a_i$, $\forall 0 \leq i \leq n$, dacché
    ogni elemento di $\FFpp$ è radice di $x^p-x$.

\end{proof}

\begin{theorem}
    Sia $f(x)$ un irriducibile di grado $n$ su $\FFpp[x]$ e sia $\alpha$ una
    sua radice in $\FFpn$. Allora vale la seguente fattorizzazione
    in $\FFpn$:

    \[ f(x) = \prod_{i=0}^{n-1} \left(x - \alpha^{p^i}\right) = \prod_{i=0}^{n-1} \left(x - \Frobexp^i(\alpha)\right), \]

    \vskip 0.1in

    dove ogni fattore non è associato.
\end{theorem}

\begin{proof}
    Si verifica innanzitutto che vale chiaramente che $\alpha^{p^i} = \Frobexp^i(\alpha)$.
    Dal momento che $\alpha$ è radice, allora ogni $\alpha^{p^i}$ lo è, per il
    \lemref{lem:frobexp}. \\

    Affinché tutti i fattori della moltiplicazione non siano associati è sufficiente
    dimostrare che $n$ è il più piccolo esponente $j$ per cui $\Frobexp^j(\alpha)=\alpha$.
    Infatti, siano $\Frobexp^i(\alpha)=\Frobexp^j(\alpha)$ con $0\leq j < i < n$, allora,
    applicando più volte $\Frob$, si ricava che:

    \[ \Frobexp^n(\alpha)=\Frobexp^{j+n-i}(\alpha) \implies \Frobexp^{j+n-i}(\alpha)=
        \alpha, \]

    \vskip 0.1in

    che è assurdo, dacché $j < i < n \implies j+n-i < n$, \Lightning{}. \\

    Innanzitutto, si verifica che $\Frobexp^{n}(\alpha)=\alpha^{p^n}=\alpha$, dacché
    $\alpha \in \FFpn$. Infine, sia $t$ il più piccolo esponente $j$ per cui
    $\Frobexp^j(\alpha)=\alpha$. Se $j$ fosse minore di $n$, $\alpha$ sarebbe
    radice di $x^{p^t}-x$. Tuttavia questo è assurdo, dal momento che così
    $\alpha$ apparterrebbe a $\FFp{t} \neq \FFpn$, quando invece il più
    piccolo campo che lo contiene è $\FFpp(\alpha) \cong \FFpp[x]/(f(x)) \cong \FFpn$,
    \Lightning{}.
\end{proof}

\subsection{L'inclusione \texorpdfstring{$\FFpm \subseteq \FFpn$}{F\_(p\string^m) in F\_(p\string^n)} e il polinomio \texorpdfstring{$x^{p^n}-x$}{x\string^(p\string^n)-x}}

\begin{lemma}
    \label{lem:alpha_radice}
    Sia $\alpha$ una radice di $x^{p^d}-x$ con $d \mid n$. Allora
    $\alpha$ è anche una radice di $x^{p^n}-x$.
\end{lemma}

\begin{proof} Sia $s \in \NN$ tale che $n=ds$.
    Si verifica la tesi applicando il principio di induzione su $k \in \NN$. \\

    \ (\textit{passo base})\; Per ipotesi, $\alpha^{p^d}=\alpha$. \\

    \ (\textit{passo induttivo})\; Per ipotesi induttiva, $\alpha^{p^{(k-1)d}}=\alpha$. Allora si ricava che:

    \[ \alpha^{p^{(k-1)d}}=\alpha \implies \alpha^{p^{kd}}=\alpha^{p^d}=\alpha. \]

    \vskip 0.1in

    In particolare, $\alpha^{p^n} = \alpha^{p^{ds}} = \alpha$, da cui la tesi.

\end{proof}

\begin{theorem}
    \label{th:inclusione}
    $\FFpm \subseteq \FFpn$ se e solo se $m \mid n$.
\end{theorem}

\begin{proof}
    Si dimostrano le due implicazioni separatamente. \\

    \ ($\implies$)\; Dal momento che $\FFpm \subseteq \FFpn$,
    si ricava la seguente catena di estensioni:

    \[ \FFpp \subseteq \FFpm \subseteq \FFpn, \]

    \vskip 0.1in

    dalla quale, applicando il \textit{Teorema delle Torri Algebriche},
    si desume la seguente equazione:

    \[ \underbrace{[\FFpn : \FFpp]}_n = [\FFpn : \FFpm] \underbrace{[\FFpm : \FFpp]}_d, \]

    e quindi che $m$ divide $n$. \\

    \ ($\,\Longleftarrow\,\,$)\; Sia $m \mid n$. Si consideri $\alpha \in \FFpm$. $\alpha$
    è sicuramente radice di $x^{p^m}-x$, e poiché $m$ divide $n$, è
    anche radice di $x^{p^n}-x$, per il \lemref{lem:alpha_radice}. Allora
    $\alpha$ appartiene al campo di spezzamento di $x^{p^n}-x$ su $\FFpp$,
    ossia $\FFpn$. Pertanto $\FFpm \subseteq \FFpn$. \\
\end{proof}

\begin{corollary}
    $\forall 1 \leq i \leq n$. Allora, detta $m_i$ il grado di $g_i(x)$, il
    campo di spezzamento di $f(x)$ è $\FFp{k}$, dove $k = \mcm(m_1, m_2, \ldots, m_n)$.
\end{corollary}

\begin{proof}
    Il campo di spezzamento di $f(x)$ è il più piccolo campo rispetto all'inclusione
    che ne contenga tutte le radici, ossia il più piccolo campo che contenga
    $\FFp{m_1}$, $\FFp{m_2}$, $\ldots,\, \FFp{m_n}$. Si dimostra che tale campo
    è proprio $\FFp{k}$. \\

    Innanzitutto $\FFp{k}$, per il \thref{th:inclusione}, contiene tutti i campi di spezzamento dei fattori irriducibili di $f(x)$, dacché $m_i$ divide $k$ $\forall 1 \leq i \leq n$. \\

    Sia supponga esista adesso un altro campo $\FFp{t} \subseteq \FFp{k}$ con tutte le
    radici. Sicuramente $t \mid k$, per il \thref{th:inclusione}. Inoltre, dal momento
    che dovrebbe includere ogni campo $\FFp{m_i}$, sempre per il \thref{th:inclusione},
    $m_i$ divide $t$ $\forall 1 \leq i \leq n$. \\

    Allora $t$ è un multiplo comune di tutti i $m_i$, e quindi $k$, in quanto minimo
    comune multiplo, lo divide. Si conclude allora che $t = k$, e quindi che
    $\FFp{k}$ è un campo di spezzamento di $f(x)$.
\end{proof}

\begin{theorem}
    $x^{p^n}-x$ è il prodotto di tutti i polinomi irriducibili in $\FFpp$
    di grado divisore di $n$.
\end{theorem}

\begin{proof}
    La proposizione è equivalente a affermare che ogni polinomio irriducibile in $\FFpp$
    ha grado divisore di $n$ se e solo se divide $x^{p^n}-x$. Si dimostrano le
    due implicazioni separatamente. \\

    \ ($\implies$)\; Sia $f(x)$ un polinomio irriducibile in $\FFpp$ di grado $d$, con
    $d \mid n$. Si consideri allora il campo $\FFpd \cong \FFpp/(f(x))$, e
    sia $\alpha$ una radice di $f(x)$ in tale campo. \\

    Per il \lemref{lem:alpha_radice} si verifica che $\alpha$ è anche una radice
    di $x^{p^n}-x$. Poiché $f(x)$ è irriducibile, esso è il polinomio minimo
    di $\alpha$, e quindi si deduce che $f(x)$ divide $x^{p^n}-x$. \\

    \ ($\,\Longleftarrow\,\,$)\; Sia $f(x)$ un polinomio irriducibile in $\FFpp$ di grado
    $d$ che divide $x^{p^n}-x$. Si consideri allora il campo $\FFpd \cong \FFpp/(f(x))$,
    e sia $\alpha$ una radice di $f(x)$ in tale campo. Allora $\FFpd \cong
        \FFpp(\alpha)$, dacché $f(x)$, in quanto irriducibile, è il polinomio minimo
    di $\alpha$. \\

    Dacché $f(x)$ divide $x^{p^n}-x$, $\alpha$ è anche una radice
    di $x^{p^n}-x$, e quindi che $\alpha \in \FFpn$. Dal momento che chiaramente
    anche $\FFpp \subseteq \FFpn$, si deduce che $\FFpd \cong \FFpp(\alpha) \subseteq
        \FFpn$. Allora, per il \thref{th:inclusione}, $d$ divide $n$.
\end{proof}

\newpage
\thispagestyle{empty}
~\newpage

\section{Polinomi simmetrici}

\subsection{Definizione e prime proprietà}

Sia $\KK$ un campo. Dati $\sigma \in S_n$ e un polinomio $f \in \KK[x_1, \ldots, x_n]$,
si definisce il seguente polinomio:

\[ (\sigma \cdot f) (x_1, \ldots, x_n) = f(x_{\sigma(1)}, \ldots, x_{\sigma(n)}), \]

\vskip 0.1in

ossia il polinomio ottenuto permutando le variabili $x_i$ secondo $\sigma$.

\begin{definition}
    Si definisce $\Sym[X_n]$ su $K$ come il sottoanello di $\KK[x_1, \ldots, x_n]$ dei
    \textbf{polinomi simmetrici}, ossia di quei polinomi tali che
    $\sigma \cdot f = f$, $\forall \sigma \in S_n$.
\end{definition}

\begin{definition}
    Sia $d \in \NN$ tale che $0 \leq d \leq n$. Si definisce \textbf{polinomio simmetrico elementare} su $\Sym[X_n]$ ogni polinomio
    della seguente forma:
    
    \[ e_d(x_1, \ldots, x_n) = \sum_{1 \leq i_1 < \cdots < i_d \leq n} \underbrace{x_{i_1} \cdots x_{i_n}}_{d\text{ volte}}, \]
    
    \vskip 0.1in
    
    dove si pone $e_0(x_1, \ldots, x_n) := 1$
\end{definition}

\begin{remark*}
    Qualora siano noti al contesto le variabili su cui è definito $\Sym[X_n]$ si
    può omettere la parentesi di $e_d$, scrivendo pertanto semplicemente
    $e_d$.
\end{remark*}

\begin{remark*}
    Sia $p(x) = a_n x^n + \ldots + a_0$ un polinomio in $\KK[x]$. Siano $\lambda_1$,
    ..., $\lambda_n$ le sue radici nel suo campo di spezzamento. Allora vale
    che:
    
    \[ a_{n-i} = (-1)^i \, a_n \, e_i(\lambda_1, \ldots, \lambda_n). \]
\end{remark*}

\begin{definition}
    Sia $\alpha = (\alpha_1, \ldots, \alpha_n) \in \NN^n$, si definisce:
    
    \[ x^\alpha = x_1^{\alpha_1} x_2^{\alpha_2} \cdots x_n^{\alpha_n}, \quad
        \card{\alpha} = \sum_{i=1}^n \alpha_i. \]
\end{definition}

\begin{remark*}
    Ogni monomio nelle variabili $x_1$, ..., $x_n$ può essere rappresentato
    nella forma $x^\alpha$, ponendo $\alpha_i$ uguale al numero di volte
    in cui la variabile $x_i$ compare nel monomio.
\end{remark*}

\begin{definition}
    Si definisce \textit{degree lexicographic order} (\textbf{deglex}) la seguente
    relazione di ordine sui monomi monici di un polinomio:

    \[ x^\alpha > x^\beta \defiff \card{\alpha} > \card{\beta} \text{ oppure }  \\
        \card{\alpha} = \card{\beta} \text{ e } \alpha > \beta \text{ secondo il LO,} \]
        
    \vskip 0.1in
    
    dove con LO si indica il \textit{lexicographic order}.
\end{definition}

\begin{proposition}
    Il \textit{deglex} è una relazione di ordine totale.
\end{proposition}

\begin{proof}
    [TODO]
\end{proof}

\begin{proposition}
    \label{prop:moltiplicazione_disuguaglianza_deglex}
    Vale la seguente equivalenza:
    
    \[ x^\alpha x^\gamma > x^\beta x^\gamma \iff x^\alpha > x^\beta. \]
\end{proposition}

\begin{proof}
    Si dimostrano le due implicazioni separamente. \\
    
    \ ($\implies$)\; Se $\card{\alpha} + \card{\gamma} > \card{\beta} + \card{\gamma}$, allora
    anche $\card{\alpha} > \card{\beta}$, e dunque
    $x^\alpha > x^\beta$. Altrimenti, esiste un $i \in \NN$ tale
    per cui $\alpha_i + \gamma_i > \beta_i + \gamma_i$ e $\alpha_j + \gamma_j
    = \beta_j + \gamma_j$ $\forall j < i$. Allora
    anche $\alpha_j = \beta_j$ $\forall j < i$ e
    $\alpha_i > \beta_i$. Dunque, per il LO, $\alpha
    > \beta$, e quindi $x^\alpha > x^\beta$. \\
    
    \ ($\,\,\Longleftarrow\,\;$)\; Se $\card{\alpha} > \card{\beta}$, allora
    anche $\card{\alpha} + \card{\gamma} > \card{\beta} + \card{\gamma}$, e dunque
    $x^\alpha x^\gamma > x^\beta x^\gamma$. Altrimenti, esiste un $i \in \NN$ tale
    per cui $\alpha_i > \beta_i$ e $\alpha_j = \beta_j$ $\forall j < i$. Allora
    anche $\alpha_j + \gamma_j = \beta_j + \gamma_j$ $\forall j < i$ e
    $\alpha_i + \gamma_i > \beta_i + \gamma_i$. Dunque, per il LO, $\alpha + \gamma
    > \beta + \gamma$, e quindi $x^\alpha x^\gamma > x^\beta x^\gamma$.
\end{proof}

\begin{proposition}
    \label{prop:numero_finito_soluzioni_deglex}

    Sia $\alpha \in \NN^n$. Allora esiste un numero finito di $\beta \in \NN^n$
    tale che $x^\alpha > x^\beta$.
\end{proposition}

\begin{proof}
    Siano fissati gli $\alpha_i$. Se $x^\alpha > x^\beta$,
    allora vale sicuramente l'equazione:
    
    \[ \alpha_1 + \ldots + \alpha_n > \beta_1 + \ldots + \beta_n, \]
    
    \vskip 0.1in
    
    che ammette un numero finito di soluzioni.
\end{proof}

\begin{definition}
    Si definisce \textbf{leading term} di un polinomio in
    $x_1$, ..., $x_n$ il termine $cx^\alpha$ tale che
    $x^\alpha > x^\beta$, per ogni altro monomio $x^\beta$
    del polinomio.
\end{definition}

\begin{proposition}
    \label{prop:leading_term_prodotto}

    Siano $f$ e $g \in \KK[x_1, \ldots, x_n]$.
    Il \textit{leading term} di $fg$ è il
    prodotto dei \textit{leading term} di $f$ e di $g$.
\end{proposition}

\begin{proof}
    Siano $x^\alpha$ e $x^\beta$ i rispettivi \textit{leading term}
    di $f$ e di $g$. Sia inoltre $x^\gamma$ il \textit{leading term}
    di $fg$. Si assuma che $x^\gamma \neq x^\alpha x^\beta$. \\
    
    Poiché ogni monomio del prodotto di $fg$ è un prodotto di due
    monomi di $f$ e di $g$, $x^\gamma$ potrà scriversi come
    prodotto di $x^\delta x^\zeta$, dove $x^\delta$ è un monomio
    di $f$ e $x^\zeta$ è un monomio di $g$. \\
    
    Poiché $x^\alpha$ è il \textit{leading term} di $f$, vale
    la seguente disuguaglianza:
    
    \[ x^\alpha > x^\delta, \]
    
    \vskip 0.1in
    
    da cui, dalla \propref{prop:moltiplicazione_disuguaglianza_deglex}, si
    ricava che:
    
    \[ x^\alpha x^\zeta > x^\delta x^\zeta. \]
    
    \vskip 0.1in
    
    Analogamente vale la seguente altra disuguaglianza:
    
    \[ x^\beta > x^\zeta, \]
    
    \vskip 0.1in
    
    da cui si ottiene che:
    
    \[ x^\alpha x^\beta > x^\alpha x^\zeta. \]
    
    \vskip 0.1in
    
    Combinando le due disuguaglianze si ottiene infine che:
    
    \[ x^\alpha x^\beta > x^\delta x^\zeta, \]
    
    \vskip 0.1in
    
    che è assurdo, dal momento che $x^\delta x^\zeta = x^\gamma$ è il \textit{leading
    term} di $fg$, \Lightning{}. Quindi $x^\gamma = x^\alpha x^\beta$.
\end{proof}

\begin{lemma}
    \label{lem:leading_term_simmetrico_disuguaglianza}
    Sia $c x^\alpha$ il \textit{leading term}
    di $f \in \Sym[X_n]$, con $\alpha = (\alpha_1, \ldots, \alpha_n) \in \NN^n$.
    Allora $\alpha_1 \geq \alpha_2 \geq \cdots \geq \alpha_n$.
\end{lemma}

\begin{proof}
    Si dimostra la tesi contronominalmente. \\
    
    Sia $c x^\beta$ un monomio di $f$ con $\beta = (\beta_1, \ldots, \beta_n)$ tale che esistano $i < j
    \mid \beta_i < \beta_j$. Si consideri $\gamma \in \NN^n$ come
    la tupla riordinata in modo decrescente di $\beta$ e sia
    $\sigma \in S_n$ tale che $\gamma = (\beta_{\sigma(1)},
    \ldots, \beta_{\sigma(n)})$. \\
    
    Poiché $f$ è un polinomio simmetrico, $\sigma \cdot f = f$. Quindi
    $f$ ammette un monomio della forma $c x^\gamma$. Dal momento
    che $\gamma > \beta$ per il LO, $x^\gamma > x^\beta$. Quindi
    $c x^\beta$ non è il \textit{leading term} di $f$.
\end{proof}

\begin{theorem}[\textit{Teorema fondamentale dei polinomi simmetrici}]
    Sia $\KK$ un campo. Vale il seguente isomorfismo:
    
    \[ \Sym[X_n] \cong \KK\left[e_1, \ldots, e_n\right]. \]
\end{theorem}

\begin{proof}
    Sia $c x^\alpha$ il \textit{leading term} di $f$, con
    $\alpha = (\alpha_1, \ldots, \alpha_n) \in \NN^n$.
    Per il \lemref{lem:leading_term_simmetrico_disuguaglianza},
    $\alpha_i - \alpha_{i+1} \geq 0$ $\forall 1 \leq i < n$. \\
    
    Si definisca dunque $\beta \in \NN^n$ in modo tale
    che $\beta_i = \alpha_i - \alpha_{i+1} \geq 0$ $\forall 1 \leq i < n$
    e $\beta_n = \alpha_n$. \\
    
    Si consideri il monomio $e_1^{\beta_1} e_2^{\beta_2} \ldots e_n^{\beta_n}$:
    il suo \textit{leading term}, per la \propref{prop:leading_term_prodotto},
    è il prodotto dei \textit{leading term} dei suoi fattori,
    ossia $x_1^{\alpha_1} \cdots x_n^{\alpha_n} = x^\alpha$. \\
    
    Si consideri adesso come polinomio $f - c e_1^{\beta_1} e_2^{\beta_2} \ldots e_n^{\beta_n}$,
    e si reiteri l'algoritmo fino a quando il risultato non è zero. Che l'algoritmo
    termini è garantito dalla \propref{prop:numero_finito_soluzioni_deglex}, da cui
    si desume che vi è numero finito di \textit{leading term} possibili una
    volta tolto ad ogni iterazione il termine $c e_1^{\beta_1} e_2^{\beta_2} \ldots e_n^{\beta_n}$. \\
    
    Infine si sarà ottenuto una rappresentazione di $f$ come combinazione di
    $e_1$, ..., $e_n$. Questa rappresentazione è unica perché
    i termini $e_1^{\beta_1} e_2^{\beta_2} \ldots e_n^{\beta_n}$ sono
    linearmente indipendenti, dal momento che i loro
    \textit{leading term} sono distinti. \\
    
    Si costruisca dunque l'omomorfismo $\Pi : \Sym[X_n] \to \KK\left[e_1, \ldots, e_n\right]$
    che associa ad ogni polinomio simmetrico la sua rappresentazione in
    $\KK\left[e_1, \ldots, e_n\right]$. \\
    
    Si verifica che $\Pi$ è un omomorfismo. Poiché tale omomorfismo è iniettivo
    e surgettivo, è un isomorfismo, da cui la tesi.
\end{proof}

\subsection{Teorema fondamentale dell'Algebra}




\newpage
\thispagestyle{empty}
~\newpage

\chapter{Riferimenti bibliografici}
\printbibliography[heading=none]

\end{document}
